% concatenated from document.tex
\documentclass[journal]{IEEEtran}

\usepackage[ruled]{algorithm2e}
\usepackage{algorithmic}
\renewcommand{\algorithmicrequire}{\textbf{Input:}}    
\renewcommand{\algorithmicensure}{\textbf{Output:}}

\usepackage{enumitem}
\usepackage[hidelinks]{hyperref}
\usepackage{graphicx}
\graphicspath{{figures/}}
\usepackage{booktabs}
\usepackage{fontawesome} 
\usepackage{multirow}
\usepackage{subcaption}

\usepackage{amsmath}
\usepackage{cases}
\usepackage{amssymb}
\usepackage{amsfonts}
\usepackage{mathtools}

\usepackage{amsthm}
\theoremstyle{definition}
\newtheorem{lemma}{Lemma}
\newtheorem{theorem}{Theorem}
\newtheorem{proposition}{Proposition}
\newtheorem{conjecture}{Conjecture}
\newtheorem{corollary}{Corollary}
\theoremstyle{definition}
\newtheorem{definition}{Definition}
\newtheorem{assumption}{Assumption}
\newtheorem{problem}{Problem}
\newtheorem*{remark}{Remark}
\newtheorem{example}{Example}

\usepackage{pifont}
\usepackage{makecell}
\newcommand{\symbox}[1]{\makebox[0.8em][c]{#1}}
\newcommand{\cmark}{\symbox{\ding{51}}}   
\newcommand{\xmark}{\symbox{\ding{55}}}   

\let\oldremark\remark
\let\endoldremark\endremark
\renewenvironment{remark}
  {\oldremark}
  {\hfill$\Diamond$\endoldremark}

\newcommand\norm[2][1]{\left\|#2\right\|_{#1}}
\newcommand\abs[1]{\left|#1\right|}
\newcommand\where[2]{\left.#1\right|_{#2}}
\newcommand\difffrac[3][1]{
	\ifnum #1=1
		\frac{\mathrm{d} #2}{\mathrm{d} #3}
	\else
	\frac{{\mathrm{d}}^{#1} #2}{\mathrm{d} #3^{#1}}
	\fi
}

\newcommand\R{{\mathbb{R}}}
\newcommand\N{{\mathbb{N}}}
\newcommand\Z{{\mathbb{Z}}}
\newcommand{\Setminus}[2]{{\left.#1\middle\backslash #2\right.}}
\newcommand\hamiltonian{{\mathcal{H}}}
\newcommand\lagrangian{{\mathcal{L}}}
\newcommand\arc{{\mathfrak{r}}}
\newcommand\singular{{\mathfrak{s}}}
\newcommand\st{{\text{s.t.}}}

\renewcommand\vector[1]{\boldsymbol{#1}}
\newcommand\va{{\vector{a}}}
\newcommand\vb{{\vector{b}}}
\newcommand\vc{{\vector{c}}}
\newcommand\vd{{\vector{d}}}
\newcommand\ve{{\vector{e}}}
\newcommand\vf{{\vector{f}}}
\newcommand\vg{{\vector{g}}}
\newcommand\vh{{\vector{h}}}
\newcommand\vm{{\vector{m}}}
\newcommand\vn{{\vector{n}}}
\newcommand\vp{{\vector{p}}}
\newcommand\vq{{\vector{q}}}
\newcommand\vr{{\vector{r}}}
\newcommand\vs{{\vector{s}}}
\newcommand\vt{{\vector{t}}}
\newcommand\vu{{\vector{u}}}
\newcommand\vw{{\vector{w}}}
\newcommand\vx{{\vector{x}}}
\newcommand\vy{{\vector{y}}}
\newcommand\vz{{\vector{z}}}
\newcommand\vA{{\vector{A}}}
\newcommand\vB{{\vector{B}}}
\newcommand\vC{{\vector{C}}}
\newcommand\vD{{\vector{D}}}
\newcommand\vE{{\vector{E}}}
\newcommand\vF{{\vector{F}}}
\newcommand\vG{{\vector{G}}}
\newcommand\vH{{\vector{H}}}
\newcommand\vI{{\vector{I}}}
\newcommand\vJ{{\vector{J}}}
\newcommand\vM{{\vector{M}}}
\newcommand\vT{{\vector{T}}}
\newcommand\vphi{{\vector{\phi}}}
\newcommand\vvarphi{{\vector{\varphi}}}
\newcommand\vlambda{{\vector{\lambda}}}
\newcommand\veta{{\vector{\eta}}}
\newcommand\vdelta{{\vector{\delta}}}
\newcommand\vzeta{{\vector{\zeta}}}
\newcommand\vxi{{\vector{\xi}}}
\newcommand\vtau{{\vector{\tau}}}
\newcommand\vone{{\vector{1}}}
\newcommand\vzero{{\vector{0}}}

\newcommand\mC{{\mathbf{C}}}
\newcommand\mG{{\mathbf{G}}}
\newcommand\mF{{\mathbf{F}}}
\newcommand\mM{{\mathbf{M}}}
\newcommand\mQ{{\mathbf{Q}}}
\newcommand\bx{{\mathbf{x}}}
\newcommand\bu{{\mathbf{u}}}

\newcommand\tf{{t_\text{f}}}
\newcommand\f{{\text{f}}}
\newcommand\s{{\text{s}}}
\newcommand\vxf{{\vx_\text{f}}}
\newcommand\sgn{{\text{sgn}}}
\newcommand\dom{{\text{dom\,}}}
\newcommand\const{{\text{const}}}

\DeclareMathOperator*{\argmax}{argmax}
\DeclareMathOperator*{\argmin}{argmin}
\DeclareMathOperator*{\minmax}{min/max}

\usepackage{accents}
\DeclareRobustCommand{\underbar}[1]{\kern0.3\fontdimen2\font\underaccent{\bar}{\kern-0.3\fontdimen2\font #1}}

\usepackage{xcolor}
\newcommand\modify[1]{\textcolor{red}{#1}}

\newcommand\RomanNum[1]{\uppercase\expandafter{\romannumeral #1}} 

\hyphenation{op-tical net-works semi-conduc-tor}

\begin{document}

\title{Reachability-Augmented Dual Dynamic Programming for Optimal Path Parameterization}

\author{
    Yunan~Wang,
    Jizhou~Yan,
    Chuxiong~Hu,
    and~Zeyang~Li
\thanks{Corresponding author: Chuxiong Hu (e-mail: cxhu@tsinghua.edu.cn).}
}

\maketitle

\begin{abstract}
    Optimal path parameterization (OPP) is a fundamental problem for planning trajectories along a prescribed geometric path under kinodynamic constraints and task-dependent objectives. While time-optimal path parameterization (TOPP) minimizes traversal time, its saturating states and controls may induce vibration and tracking errors, which can be mitigated by introducing smoothness objectives. However, a key capability gap remains in OPP: feasibility guarantees, general-objective optimality certificates, and computational efficiency are difficult to achieve simultaneously in a unified framework, especially for third-order OPP (OPP3) with non-convex constraints. This paper proposes reachability-augmented dual dynamic programming (RDDP), a state-grid-free and objective-aware dynamic programming (DP) framework for OPP. The key idea is to replace the relatively complete recourse assumption used in classical dual DP (DDP) with OPP-specific backward reachable sets, and then generate both value-function cuts and trial trajectories only inside these reachable sets. For convex and non-convex OPP, we prove global optimality and Karush-Kuhn-Tucker convergence of RDDP under OPP-specific conditions, respectively. Efficient instantiations are developed for second-order OPP (OPP2) and OPP3. Experiments show that RDDP achieves objective values comparable to convex-optimization baselines while reducing computation time by 28.6 times for OPP2 and 5.8 times for OPP3. RDDP also achieves faster convergence than grid-based DP. Compared with reachability-analysis methods, RDDP retains the reachability mechanism while replacing local maximum-control propagation with value-function-guided control selection, thereby enabling objectives beyond traversal time. In summary, RDDP addresses a key capability gap in OPP by unifying certifiable general-objective optimization, reachability-based feasibility preservation, and online-compatible low-dimensional DP computation in a single OPP framework.
\end{abstract}

\begin{IEEEkeywords}
    Optimal control, path parameterization, dual dynamic programming, reachability analysis, general objective.
\end{IEEEkeywords}

\IEEEpeerreviewmaketitle

\section{Introduction}\label{sec:Introduction}

\IEEEPARstart{O}{ptimal} path parameterization (OPP) is a fundamental problem in robotics and optimal control. It aims to find an optimal time parameterization of a given geometric path under kinodynamic constraints and task-dependent objectives. Originating from the study in the 1980s \cite{bobrow1985time,shin1985minimum,shin1986dynamic}, OPP has been widely applied in various robotic systems, such as manipulators \cite{verscheure2009time,debrouwere2013time,wang2022third}, humanoid robots \cite{pham2018new,hauser2014fast}, and manufacturing systems \cite{wang2026online,xu2025real}.

When the objective is to minimize the traversal time, OPP reduces to the well-known time-optimal path parameterization (TOPP) problem \cite{pham2018new}. According to the Pontryagin's maximum principle (PMP) \cite{pontryagin1987mathematical}, the optimal control of TOPP is bang-singular-bang, which pushes the velocity, acceleration, torque, and jerk constraints to their limits. Despite the high motion efficiency, such constraint-saturating trajectories might induce undesirable behaviors such as vibration and tracking error in some scenarios. By trading traversal time with task-dependent criteria such as energy consumption and smoothness \cite{verscheure2009time,kaserer2019nearly}, general objectives can smooth the motion profile while maintaining high motion efficiency. Our real-world experiment provides one example of this tradeoff: compared with pure TOPP, the time-thermal objective increases the traversal time by only 0.15\% while reducing the normalized total variation of feedback acceleration by 87.7\%.

For online robotic applications, an ideal OPP solver should satisfy four requirements simultaneously:
\begin{enumerate}
    \item \textbf{Feasibility}: The solver should guarantee the feasibility of the generated trajectory under state and mixed state-control constraints, which are essential for the safety and reliability of robotic systems.
    \item \textbf{General-objective optimality}: The solver should provide appropriate optimality certificates for general objectives.
    \item \textbf{Efficiency}: The solver should be computationally efficient enough for online motion generation.
    \item \textbf{High-order capability}: The solver should be able to handle higher-order systems, especially third-order OPP (OPP3) which involves non-convex jerk constraints.
\end{enumerate}
To the best of the authors' knowledge, it remains challenging to simultaneously achieve all the above properties within a unified framework. This key capability gap is particularly pronounced for OPP3 due to non-convexity.

\begin{table*}[!t]
    \centering
    \caption{Comparison of different classes of OPP solvers.}
    \label{tab:related_works}
    \begin{tabular}{llllll}
        \toprule
        Method & Order & Objective Scope & Optimality & Feasibility & Efficiency \\ \midrule
        NI \cite{pham2014general} & 2 & \xmark\,(Time only) & \cmark\,(Global) & \cmark & High \\
        CO \cite{verscheure2009time,hauser2014fast} & 2 & \cmark\,(General) & \cmark\,(Global) & \cmark & Low \\
        DP \cite{shin1986dynamic} & 2 & \cmark\,(General) & \cmark\,(Grid-optimal) & \cmark & Medium \\
        RA \cite{pham2018new} & 2 & \xmark\,(Time only) & \cmark\,(Global, without zero-inertia points) & \cmark & High \\
        \textbf{RDDP (Ours)} & 2 & \cmark\,(General) & \cmark\,(Global) & \cmark & Medium--High \\ \midrule
        NI \cite{pham2017structure} & 3 & \xmark\,(Time only) & \xmark & \xmark & N/A \\
        CO \cite{debrouwere2013time,wang2026online} & 3 & \cmark\,(General) & \cmark\,(KKT convergence) & \cmark & Low \\
        DP \cite{shin1986dynamic} & 3 & \cmark\,(General) & \cmark\,(Grid-optimal) & \cmark & Extremely Low \\
        Greedy DP \cite{kaserer2019nearly} & 3 & \cmark\,(General) & \xmark & \xmark & Medium \\
        RA \cite{dio2025time} & 3 & \xmark\,(Time only) & \cmark\,(Empirically near optimal) & \cmark & High \\
        \textbf{RDDP (Ours)} & 3 & \cmark\,(General) & \cmark\,(KKT convergence) & \cmark & Medium \\ \bottomrule
    \end{tabular}
\end{table*}

\subsection{Related Works}\label{subsec:related_works}

\subsubsection{OPP solvers}\label{sub2sec:review_opp_solvers}

The study of OPP, including TOPP, can be categorized into three main classes: numerical integration (NI), convex optimization (CO), and dynamic programming (DP). In this paper, we classify the recently developed reachability analysis (RA) approaches as DP methods, as they involve the backward propagation of reachable sets in a DP manner.

As the earliest approaches to TOPP \cite{shin1985minimum}, NI methods exploit the bang-singular-bang structure of the optimal control induced by PMP and reduce the decision space to a binary choice between maximum- and minimum-control arcs. As a result, NI methods are highly efficient for TOPP2, where time-optimality is guaranteed by the consistent maximum property of the optimal solution. Despite numerical instability issues induced by the dynamic singularity \cite{shiller1992computation}, some robust TOPP2-NI implementations have been developed \cite{pham2014general}. However, a complete extension of NI methods to TOPP3 is still lacking \cite{pham2017structure}. NI methods also do not naturally extend to general-objective OPP, as the bang-singular-bang structure no longer holds. This reveals the \textit{objective-agnostic} nature of NI methods.

CO approaches provide the strongest generic optimality certificates among existing OPP solvers. The standard optimization formulations for second-order OPP (OPP2) and OPP3 were established in \cite{verscheure2009time} and \cite{debrouwere2013time}, respectively. These formulations are highly expressive \cite{hauser2014fast}. For convex OPP, CO methods can guarantee global optimality; for non-convex OPP, CO methods can obtain Karush-Kuhn-Tucker (KKT) solutions based on the sequential convex programming (SCP) framework \cite{le2018dc}. For online applications with long paths, recent advances have further guaranteed the feasibility of the windowing process for non-convex OPP3 \cite{wang2026online}. As a result, CO approaches satisfy three of the requirements stated above, i.e., optimality, feasibility, and high-order capability. However, the main limitation of CO methods lies in their computational efficiency, especially for OPP3 which involves non-convex third-order constraints. Although OPP is often used as a low-level motion planning module that must be repeatedly solved with short latency, the considered optimization problems can involve thousands of decision variables and constraints. Such large-scale non-convex optimization problems remain computationally expensive for online OPP applications.

The above observation motivates the study of DP methods. Instead of solving the full large-scale optimization problem in OPP, DP methods decompose it into a sequence of low-dimensional subproblems. Standard DP methods discretize the state space into grid points at each parameter point \cite{shin1986dynamic}. Based on Bellman's principle of optimality (BPO) \cite{bellman1952theory}, DP methods generate the optimal grid sequence under general objectives based on the value function. Theoretically, grid-based DP can approach the global optimum in the continuous state space as the grid resolution tends to infinity; however, the grid refinement suffers from slow optimality convergence and high computational cost. Furthermore, the state space in OPP3 is two-dimensional, thereby exhibiting curse of dimensionality. To address this issue, some works improve efficiency by greedily pruning grid states and reducing the two-dimensional grid into a one-dimensional subset \cite{kaserer2019nearly}, which sacrifices feasibility and optimality guarantees. In summary, grid-based DP methods can handle general objectives, but they cannot guarantee optimality for OPP3 within acceptable computational time due to the curse of dimensionality.

In contrast to grid-based DP methods, RA methods propagate backward reachable sets directly in the continuous state space, thereby providing a state-grid-free mechanism for feasibility preservation. For convex TOPP2, the backward reachable sets can be exactly computed and represented as intervals \cite{pham2018new,consolini2019optimal}. For non-convex TOPP3, existing RA works typically compute conservative inner approximations of the backward reachable sets using two-dimensional polytopes \cite{dio2025time}. Once backward reachable sets are computed, RA methods perform forward propagation by locally maximizing the admissible control, which is a key reason for the high efficiency of RA. Benefiting from the consistent maximum property of TOPP2, RA methods can guarantee near-optimality for TOPP2, as a larger feasible path velocity is aligned with the minimum-time objective. Although such monotonicity does not hold in TOPP3, near optimality can be empirically achieved. However, the limitation appears in general OPP when the objective is no longer purely traversal time. Therefore, the open challenge is not feasibility preservation, where RA is already effective, but how to retain RA's reachability mechanism while achieving \textit{objective-aware} control selection for general OPP.

Table \ref{tab:related_works} summarizes the above review and illustrates the capability gap addressed in this paper. The computational efficiency is compared qualitatively within each order. More quantitative details of different methods are provided in Section \ref{subsec:numerical_experiments}.

\subsubsection{Dual Dynamic Programming (DDP)}\label{sub2sec:review_dpp}

DDP \cite{read1990dual}, also known in the stochastic setting as stochastic DDP (SDDP) \cite{pereira1991multi}, is a state-grid-free DP method for multistage convex optimization. In contrast to conventional grid-based DP methods, DDP approximates the value function from below using affine optimality cuts over the continuous state space. Then, the optimality cuts are repeatedly refined around the trial trajectory generated by the current approximate value function, which partially alleviates the curse of dimensionality.

From an optimization perspective, DDP can be interpreted as nested Benders decomposition \cite{bnnobrs1962partitioning,birge1985decomposition}, which approximate convex value functions by affine cuts. Instead of solving the large-scale optimization problem as a whole, DDP decomposes it into a sequence of subproblems while preserving convergence to the globally optimal value \cite{fullner2025stochastic} under standard assumptions.

A central feasibility issue of standard DDP is the relatively complete recourse assumption \cite{rockafellar1976stochastic}, under which every feasible current state admits a feasible continuation in future stages. Once such an assumption fails, feasibility cuts are introduced to exclude states that cannot be extended feasibly to the terminal set \cite{guigues2014sddp}, where an outer approximation of recourse-feasible region is iteratively refined. This mechanism is effective for general multistage optimization problems if the recourse-feasible region is not explicitly available.

To the best of the authors' knowledge, however, DDP has not been introduced to robotic OPP so far. This absence reflects a structural challenge rather than merely a lack of attention, since directly applying the classical DDP to OPP is nontrivial. As will be analyzed in Section \ref{subsec:key_ideas}, the difficulty rises from the hard and typically degenerate terminal set, together with the non-convexity of OPP. These observations motivate an OPP-specific DDP framework in this paper.

\subsection{Challenges and Key Ideas}\label{subsec:key_ideas}

As discussed in Section \ref{sub2sec:review_opp_solvers}, existing RA approaches in TOPP \cite{pham2018new,dio2025time} provide an efficient state-grid-free mechanism for feasibility and reachability. The missing component for general-objective OPP is a cost-to-go model. This motivates a combination of the value function with the RA framework, thereby preserving the feasibility guarantees of conventional RA while additionally enabling objective-aware control selection. Once the exact value function is available for every state in the backward reachable set, forward propagation can select the optimal admissible control at each stage, which results in the global optimality by BPO \cite{bellman1952theory}. However, the exact value function is computationally intractable in practice, as the backward reachable set contains uncountably many states.

To address the above issue, a state-grid-free DP framework should be introduced to mitigate the curse of dimensionality. As reviewed in Section \ref{sub2sec:review_dpp}, DDP is a potential candidate for this purpose, as it converges to the global optimum for convex multistage problems under suitable assumptions. However, the direct application of classical DDP to OPP is nontrivial due to feasibility and non-convexity issues.

The first difficulty lies in recourse-feasibility which is closely related to reachability in OPP. The terminal set is typically a singleton in OPP, i.e., the terminal state is fully specified. In principle, the lack of recourse in DDP can be represented by an extended-real value function, which assigns $+\infty$ to states outside the backward reachable set. However, the extended-real value function is infinite almost everywhere near the terminal stage due to the low-dimensionality of the backward reachable set, which leads to numerical failures in guiding the forward trajectory generation. As an alternative, feasibility cuts can be introduced to avoid the reliance on the relatively complete recourse assumption in OPP, where the approximation of backward reachable sets should be iteratively refined. Such approach is computationally expensive and may not provide feasible trajectories in early iterations, thus being undesirable for online applications. Furthermore, slack variables and penalty reformulations are also unsuitable for OPP due to the hard stage and terminal constraints.

Another difficulty rises from the non-convexity of OPP, as the standard DDP is primarily designed for convex problems, where affine optimality cuts are valid lower bounds of the convex value functions. This paper embeds DDP into the SCP framework to handle non-convex OPP. However, as will be analyzed in Section \ref{subsec:KKT_GOPP_Conservative}, the global optimality of each convexified subproblem might be unavailable in an RA-based DDP framework. In this case, the optimality convergence of SCP therefore requires a careful analysis.

In summary, RA and DDP address complementary components of the capability gap in general-objective OPP. RA provides explicit reachable sets for feasibility preservation, but it is specialized to time objectives. DDP provides objective-aware optimality guarantees for convex problems, but its classical formulation does not directly handle the reachability-critical and non-convex structure of OPP. These observations motivate the development of a reachability-augmented DDP framework in this paper, which combines the explicit feasibility preservation of RA with the objective-aware optimality of DDP. The resulting framework unifies reachability-based feasibility preservation, certifiable general-objective optimization, and online-compatible low-dimensional DP computation for both OPP2 and OPP3.

\subsection{Contributions and Comparison with Existing Methods}

The comparison of the proposed framework with existing methods is summarized in Table \ref{tab:related_works}. The main contributions and comparison with existing works are as follows:
\begin{enumerate}
    \item \textbf{Novel OPP framework augmenting DDP with RA.} To the best of the authors' knowledge, this paper introduces a DDP formulation to robotic OPP for the first time. To avoid relying on recourse assumptions or feasibility cuts in DDP that are unsuitable for OPP, this paper develops an OPP-specific reachability-augmented DDP (RDDP) framework. The key difference from a direct application of DDP is that RDDP explicitly uses backward reachable sets as recourse-feasible domains, inside which value-function cuts and trial trajectories are generated. Compared with RA, RDDP preserves reachability but replaces local maximum-control propagation with value-function-guided control selection, thereby enabling general objectives beyond traversal time. Compared with grid-based DP, the state-grid-free RDDP does not require the generated trajectory to pass through state-grid points. Compared with CO, RDDP avoids solving the full large-scale optimization problem by decomposing OPP into a sequence of low-dimensional subproblems. Therefore, RDDP unifies certifiable general-objective optimization, reachability-based feasibility preservation, and online-compatible low-dimensional DP computation in a single state-grid-free framework for OPP.
    \item \textbf{Feasibility and optimality guarantees.} Since DDP replaces the standard recourse assumption of DDP with explicit reachability information, this paper establishes theoretical feasibility and optimality guarantees for RDDP under new OPP-specific conditions. For convex OPP (COPP), the historical best solution is proved to converge to the global optimum of COPP. Benefiting from the introduced reachable set, each forward iteration of RDDP generates a feasible trajectory with computable DDP optimality gap. This property is particularly important for online OPP, as it enables practical early termination under limited computational budgets. Note that feasibility guarantee is generally not provided in a direct DDP implementation before sufficiently many feasibility cuts have been generated. For general non-convex OPP (GOPP), RDDP is embedded into an SCP framework with feasibility guarantee and optimality convergence to KKT solutions. The analysis is further extended to conservative approximate reachable sets, which are practically important for OPP3. Therefore, RDDP achieves comparable optimality level to CO methods, i.e., global optimality and KKT convergence in convex and non-convex OPP, respectively.
    \item \textbf{Efficient instantiations for OPP2 and OPP3.} This paper develops concrete RDDP algorithms for the widely studied OPP2 and OPP3. For OPP2, exact interval reachable sets and value-function cuts are computed through one- and two-dimensional convex optimization, respectively. For OPP3, conservative polytopic reachable sets are combined with projection-based and searching-based routines to control computational complexity, where convex optimization problems of dimension no greater than three are solved. In numerical experiments, RDDP achieves objective values comparable to CO baselines while achieving a speedup of approximately 28.6 and 5.8 times for OPP2 and OPP3, respectively. RDDP also converges faster than grid-based DP. Compared with RA, RDDP handles broader objectives for OPP. The developed RDDP can serve as a state-grid-free and objective-aware solver for general-objective OPP2 and OPP3 with feasibility and optimality for online robotic applications.
\end{enumerate}

\section{Problem Formulation and Preliminaries}\label{sec:formulation_preliminaries}

This section formulates the considered OPP problem and provides the theoretical preliminaries. Section \ref{subsec:formulation_opp} formulates OPP as an optimization problem. Then, the concepts of COPP and GOPP are formally stated in Section \ref{subsec:problem_statement}. Finally, two structural components of the proposed RDDP framework, i.e., the value function and the backward reachable set, are introduced in Section \ref{subsec:preliminary_dp}.

\subsection{Formulation of OPP}\label{subsec:formulation_opp}

Consider an $n$-dimensional robotic system with configuration $\vq\in\R^n$. A regular geometric path $\vq=\vq(s)$ of $\mathcal{C}^m$ smoothness is given, where $s\in[0,s_\text{f}]$ is the path parameter and the terminal parameter $s_\text{f}$ is fixed. Note that $s$ is not necessarily the arc length. The OPP problem is to find a strictly increasing time parameterization $s=s(t)$ that minimizes a user-specified objective under state and control constraints, where $t\in[0,t_\text{f}]$ is the time and the terminal time $t_\text{f}$ is free.

Following the standard process in \cite{verscheure2009time,debrouwere2013time}, the above description in the $t$-domain can be transformed into the $s$-domain, which is more convenient for numerical solution. The details of the transformation are provided in Appendix \ref{app:applicability}. As a result, the parameter domain $[0,s_\text{f}]$ is discretized into a grid $0=s_0<s_1<\cdots<s_N=s_\text{f}$. Denote the state at $s_k$ by $\vx(s_k)=\vx_k\in\R^{m-1}$ and the control in $(s_k,s_{k+1})$ by $u(s)\equiv u_k$. Specifically, $(\vx,u)$ is the discretization of $\Big(\frac12\dot{s}^2,\frac{\mathrm{d}}{\mathrm{d}s}\Big(\frac12\dot{s}^2\Big),\dots,\frac{\mathrm{d}^{m-1}}{\mathrm{d}s^{m-1}}\Big(\frac12\dot{s}^2\Big)\Big)$. The resulting optimization problem in the $s$-domain is given by
\begin{subequations}\label{eq:op_discretize}
    \begin{align}
        \min_{\vx_k,u_k} \quad &J = \sum_{k=0}^{N-1} L_k(\vx_k, u_k) + \Phi(\vx_N) \label{eq:op_discretize_objective}\\
        \st \quad & \vx_{k+1} = \vA_k\vx_k + \vb_k u_k, \,\,k=0,1,\ldots,N-1,\label{eq:op_discretize_dynamics} \\
        &\vf_k(\vx_k) \leq \vzero,  \,\,k=1,2,\ldots,N-1, \label{eq:op_discretize_state_constraints}\\
        &\vh_k(\vx_k, u_k) \leq \vzero, \,\,k=0,1,\ldots,N-1,\label{eq:op_discretize_mix_constraints}\\
        &\vx_0\in\mathcal{X}_0, \,\,\vx_N\in\mathcal{X}_\text{f}.
    \end{align}
\end{subequations}
According to the zero-hold of $u(s)$ in each interval, the discretized chain-of-integrator dynamics \eqref{eq:op_discretize_dynamics} is \cite{wang2025time}
\begin{equation}\label{eq:Ab}
    \vA_k=\left[\begin{array}{cccc}
        1 & \Delta_k & \cdots & \frac{\Delta_k^{m-2}}{(m-2)!} \\
        0 & 1 & \cdots & \frac{\Delta_k^{m-3}}{(m-3)!} \\
        \vdots & \vdots & \ddots & \vdots \\
        0 & 0 & \cdots & 1
    \end{array}\right],\,\,\vb_k=\left[\begin{array}{c}
        \frac{\Delta_k^{m-1}}{(m-1)!} \\
        \frac{\Delta_k^{m-2}}{(m-2)!} \\
        \vdots \\
        \Delta_k
    \end{array}\right],
\end{equation}
where $\Delta_k= s_{k+1}-s_{k}>0$. All the functions and sets in problem \eqref{eq:op_discretize} are given, whereas $\vx_k$ and $u_k$ for all $k$ are the decision variables. The applicability of formulation \eqref{eq:op_discretize} to general OPP problems is discussed in Appendix \ref{app:applicability}.

\begin{assumption}\label{assum:feasibility_smoothness}
    Functions $L_k$, $\Phi$, $\vf_k$, and $\vh_k$ are locally Lipschitz continuous in their arguments for each $k$. The sets $\mathcal{X}_0$, $\mathcal{X}_\text{f}$, and $\{(\vx_k,u_k)\mid \vf_k(\vx_k)\leq\vzero,\,\vh_k(\vx_k,u_k)\leq\vzero\}$ are compact for each $k$.
\end{assumption}

\begin{remark}
    Most OPP problems of interest satisfy Assumption \ref{assum:feasibility_smoothness}. For example, the state and control are usually bounded by kinodynamic limits. The objective and constraint functions are usually locally Lipschitz continuous.

    Two exceptions need to be clarified. Denote the derivative of a variable $\bullet$ with respect to $t$ as $\dot{\bullet}=\frac{\mathrm{d}\bullet}{\mathrm{d}t}$. First, for stationary boundary conditions in OPP3, i.e., $\dot{s}=0$ at $s=0$ and $s=s_\text{f}$, the control $u=\frac{\dddot{s}}{\dot{s}}$ can be unbounded even if the jerk limit \eqref{eq:jerk_constraint} is bounded. In this case, the zero-hold of $\frac{\dddot{s}}{\dot{s}}$ in the boundary interval can be replaced by a zero-hold of $\dddot{s}$ \cite[Section 3.2.2]{wang2026online}, where the dynamic matrix and vector in \eqref{eq:Ab} are modified accordingly. Second, for some objectives, such as time objective \eqref{eq:objective_time}, the objective function is ill-defined at the boundary of the state space. For example, $L_k(\vx_k, u_k)=\frac{1}{\sqrt{2x_{k,1}}}$ is not well-defined at $x_{k,1}=0$. In this case, we consider a lower level set of $L_k$, such as $\{x_{k,1}\geq 10^{-12}\}$, to avoid the singularity, the regularization of which is chosen small enough not to affect optimality. So the above two exceptions can be handled by minor modifications of formulation \eqref{eq:op_discretize}, and thus do not violate Assumption \ref{assum:feasibility_smoothness}.
\end{remark}

\subsection{Problem Statement}\label{subsec:problem_statement}

This paper considers two classes of OPP \eqref{eq:op_discretize}, i.e., convex and non-convex problems. Both of the following problems are considered under Assumption \ref{assum:feasibility_smoothness}. Some common OPP2 problems are typically convex, whereas some typical OPP3 problems are non-convex due to the non-convex jerk constraints \eqref{eq:jerk_constraint}.

\begin{problem}[Convex optimal path parameterization, COPP]\label{problem:convex}
    Functions $L_k$, $\Phi$, $\vf_k$, and $\vh_k$ are convex in their arguments for each fixed $k$. Sets $\mathcal{X}_0$ and $\mathcal{X}_\text{f}$ are convex as well. Problem \eqref{eq:op_discretize} is feasible. The goal is to find a globally optimal solution to COPP \eqref{eq:op_discretize}.
\end{problem}

\begin{definition}\label{def:DC_decomposable}
    \cite{sriperumbudur2009convergence}
    Consider a $\mathcal{C}^1$-continuous and non-convex function $f:\Omega\subset\R^n\to\R$ where $\Omega$ is a non-empty convex set. If $f$ can be expressed as the difference of two convex functions (DC), then $f$ is called DC-decomposable. In other words, there exist two $\mathcal{C}^1$-continuous convex functions $f_1,f_2:\Omega\to\R$ such that $f(\vx)=f_1(\vx)-f_2(\vx)$ for all $\vx\in\Omega$. Given $\vx^*\in\Omega$, we define the convexified function of $f$ at $\vx^*$ as follows:
    \begin{equation}
        \overline{f}(\vx;\vx^*)=f_1(\vx)-\left[f_2(\vx^*)+\nabla f_2(\vx^*)^\top(\vx-\vx^*)\right].
    \end{equation}

    Consider a vector-valued DC-decomposable function $\vf=(f_i)_{i=1}^m:\Omega\subset\R^n\to\R^m$, i.e., each component function $f_i$ is DC-decomposable. We define the convexified function of $\vf$ at $\vx^*$ as $\overline{\vf}(\vx;\vx^*)=(\overline{f}_i(\vx;\vx^*))_{i=1}^m$.

    A set $\mathcal{X}=\{\vx\in\R^n\mid \vf(\vx)\leq\vzero,\,\vg(\vx)=\vzero\}$ is called DC-decomposable, if $\vf$ is DC-decomposable and $\vg$ is affine. We define the convexified set of $\mathcal{X}$ at $\vx^*$ as $\overline{\mathcal{X}}(\vx^*)=\{\vx\in\R^n\mid \overline{\vf}(\vx;\vx^*)\leq\vzero,\,\vg(\vx)=\vzero\}$.
\end{definition}

\begin{problem}[General optimal path parameterization, GOPP]\label{problem:general}
    Functions $L_k$, $\Phi$, $\vf_k$, and $\vh_k$ are DC-decomposable. Sets $\mathcal{X}_0$ and $\mathcal{X}_\text{f}$ are also DC-decomposable. The goal is to find a KKT solution to GOPP \eqref{eq:op_discretize}.
\end{problem}

\begin{remark}
    The smoothness assumption in GOPP is stronger than that in COPP for the well-posedness of the KKT condition. In practice, such smoothness requirement is inherently satisfied by the analytical expressions of the objective and constraint functions. As for some special objectives, such as the total variation of the torque \eqref{eq:objective_torque_variation}, $L_k$ might be non-smooth at some points. In this case, $L_k$ can be approximated by a smooth variant, such as a smoothed maximum operator \cite{nesterov2005smooth} or a piecewise polynomial function.
\end{remark}

Both COPP and GOPP have been well addressed in CO-based literature \cite{verscheure2009time,debrouwere2013time,wang2026online} by leveraging off-the-shelf general-purpose solvers. As analyzed in Section \ref{sec:Introduction}, this paper focuses on developing an efficient objective-aware DP framework for solving OPP, where the large-scale optimization problem \eqref{eq:op_discretize} should be decomposed into a sequence of low-dimensional local subproblems.

\subsection{Reachable Set and Value Function}\label{subsec:preliminary_dp}

The RDDP is built on two structural components of DP \cite{bellman1952theory}, i.e., the backward reachable set and the value function. The former provides feasibility preservation, while the latter enables general-objective optimization.

\subsubsection{Backward Reachable Set} To guarantee the feasibility, the backward reachable set \cite{bokanowski2010reachability} is introduced as follows.

\begin{definition}\label{def:reachable_set}
    The backward reachable set $\mathcal{B}_k$ at $s_k$ is the set of states $\vx_k$ from which there exists a feasible trajectory that can reach the terminal set $\mathcal{X}_\text{f}$, i.e., $\forall k=0,1,\dots,N$,
    \begin{align}
        \mathcal{B}_k=\Big\{&\vx_k\in\R^{m-1}\mid\exists \vx_{k+1}, \ldots, \vx_N, u_k, \ldots, u_{N-1}, \notag\\
        &\text{ s.t. } \forall k\leq j<N,\,\vx_{j+1} = \vA_j\vx_j+\vb_j u_j, \notag\\
        &\qquad\vf_j(\vx_j) \leq \vzero,\,\vh_j(\vx_j, u_j) \leq \vzero,\,\vx_N \in \mathcal{X}_\text{f} \Big\}.\label{eq:backward_reachable_definition}
    \end{align}
    By convention, the terminal reachable set is $\mathcal{B}_N=\mathcal{X}_\text{f}$.
\end{definition}

\begin{figure}[!]
    \centering
    \includegraphics[width=0.9\columnwidth]{dp.png}
    \caption{Structural components of RDDP. (a) The backward reachable set $\mathcal{B}_k$ at $s_k$ is propagated backward from the terminal set $\mathcal{X}_\text{f}$ at $s_N$. The forward propagation based on $\mathcal{B}_k$ provides a sufficient and necessary condition for feasibility. (b) The value function $V_k^*$ is propagated backward from the terminal cost $\Phi$ at $s_N$. The forward propagation based on $V_k^*$ provides a sufficient and necessary condition for optimality.}
    \label{fig:dp}
\end{figure}

According to Definition \ref{def:reachable_set}, for $k<N$, the backward reachable set $\mathcal{B}_k$ satisfies the following backward recursion:
\begin{align}
    \mathcal{B}_k=\Big\{&\vx_k\in\R^{m-1}\mid\exists u_k,\text{ s.t. } \vh_k(\vx_k, u_k) \leq \vzero, \notag\\
    &\vf_k(\vx_k)\leq\vzero,\,\vx_{k+1}=\vA_k\vx_k+\vb_k u_k\in\mathcal{B}_{k+1}\Big\}.\label{eq:recursion_backward}
\end{align}
For each state $\vx_k \in \mathcal{B}_k$, the admissible control set is
\begin{align}
    \mathcal{C}_k(\vx_k) =\Big\{&u_k \in \R \mid \vh_k(\vx_k, u_k) \leq \vzero, \notag\\
    &\vx_{k+1}=\vA_k\vx_k+\vb_k u_k\in\mathcal{B}_{k+1}\Big\}.\label{eq:admissable_control_set}
\end{align}
Then, \eqref{eq:recursion_backward} means that $\vx_k\in\mathcal{B}_k$ if and only if $\vx_k$ satisfies the state constraints \eqref{eq:op_discretize_state_constraints} and $\mathcal{C}_k(\vx_k)\neq\varnothing$. 

In fact, \eqref{eq:recursion_backward} provides a backward propagation (BP) approach to compute the backward reachable set $\mathcal{B}_k$ for each $k$. As shown in Fig. \ref{fig:dp}(a), $\mathcal{B}_k$ can be determined by $\mathcal{B}_{k+1}$ and the constraints of $(\vx_k,u_k)$ at $s_k$. Once all $\mathcal{B}_k$ are computed, one can directly generate feasible trajectories by forward propagation (FP) from $\mathcal{B}_0$ to $\mathcal{B}_N$. Specifically, an initial state $\vx_0 \in \mathcal{B}_0\cap \mathcal{X}_0$ is first chosen, and then for each $k$, a control $u_k\in\mathcal{C}_k(\vx_k)\neq\varnothing$ is chosen to get $\vx_{k+1}=\vA_k\vx_k+\vb_k u_k\in\mathcal{B}_{k+1}$ until $s_N$. In this way, $\forall k$, the state $\vx_k$ is guaranteed to be in $\mathcal{B}_k$, and thus the generated trajectory is feasible.

\begin{remark}
    For general optimal control problems, $\mathcal{B}_k$ has no analytical expression and is computationally expensive \cite{bokanowski2010reachability}. The specific structure of second- (COPP2) and third-order COPP (COPP3) permits its efficient evaluation and representation \cite{pham2018new,dio2025time}. More details are given in Section \ref{sec:ardp_opp2_opp3}.
\end{remark}

\subsubsection{Value Function}

The value function \cite{bellman1962applied} $V_k^*(\vx_k)$ is introduced for the optimality as follows.

\begin{definition}
    The value function $V_k^*(\vx_k)$ is defined as the optimal cost-to-go from state $\vx_k$ at $s_k$, i.e., $V_k^*:\mathcal{B}_k \to \R$,
    \begin{subequations}\label{eq:def_value_function}
        \begin{align}
            V_k^*(\vx_k) = \min_{(\vx_{j+1},u_j)_{j\geq k}}&\sum_{j=k}^{N-1} L_j(\vx_j, u_j) + \Phi(\vx_N) \\
            \st\,\,&\vx_{j+1} = \vA_j\vx_j + \vb_j u_j,\,\,j\geq k,\\
            &\vf_j(\vx_j) \leq \vzero,\,\,j> k,\\
            &\vh_j(\vx_j, u_j) \leq \vzero,\,\,j\geq k,\\
            &\vx_N \in \mathcal{X}_\text{f}.
        \end{align}
    \end{subequations}
    By convention, the terminal value function $V_N(\vx_N)$ is defined as $V_N(\vx_N)=\Phi(\vx_N)$ for $\vx_N\in\mathcal{X}_\text{f}=\mathcal{B}_N$.
\end{definition}

The Bellman's optimality equation implies that for $k<N$,
\begin{equation}\label{eq:bellman_value_function}
    V_k^*(\vx_k) = \min_{u_k\in\mathcal{C}_k(\vx_k)} L_k(\vx_k, u_k) + V_{k+1}^*(\vA_k\vx_k+\vb_k u_k).
\end{equation}

As shown in Fig. \ref{fig:dp}(b), the value function $V_k^*$ can be also computed via BP recursively from $s_N$ to $s_0$ by \eqref{eq:bellman_value_function}. Once all $V_k^*$ at each state $x_k$ are computed, one can generate a globally optimal trajectory by FP from $s_0$ to $s_N$. Specifically, an initial state is chosen by $\vx_0^*\in \argmin_{\vx_0\in\mathcal{B}_0\cap \mathcal{X}_0}V^*(\vx_0)$. For each $k$, an optimal control is selected by $u_k^* \in \argmin_{u_k\in\mathcal{C}_k(\vx_k^*)}\Big[ L_k(\vx_k^*, u_k) + V_{k+1}^*(\vA_k\vx_k^*+\vb_k u_k)\Big]$ to get $\vx_{k+1}^*=\vA_k\vx_k^*+\vb_k u_k^*$ until $s_N$. The global optimality is guaranteed by the BPO \cite{bellman1962applied}.

\begin{figure}[!t]
    \centering
    \includegraphics[width=0.99\columnwidth]{ardp_one_loop.png}
    \caption{RDDP process for COPP. (a) The dependence relationship in RDDP. (b) The relation between the real value function $V_k^*$, the approximate value function $\mathcal{V}_k^p$, and the one-step approximate value function $\widehat{\mathcal{V}}_k^p$. The DDP optimality gap $\varepsilon_\text{gap}^k$ is defined in \eqref{eq:dual_gap} and \eqref{eq:dual_gap_k}.}
    \label{fig:ardp_one_loop}
\end{figure}

\section{RDDP for COPP with Global Optimality}\label{sec:RDDP_COPP}

To address the intractability of ideal DP in Section \ref{subsec:preliminary_dp}, this section proposes the RDDP framework to solve COPP, i.e., the convex version of OPP defined in Problem \ref{problem:convex}. In contrast to the classical DDP where the relatively complete recourse assumption or feasibility cuts are used to handle the feasibility, RDDP explicitly uses the backward reachable set to guarantee the feasibility of forward trajectories, as presented in Section \ref{subsec:rddp_process}. Under such OPP-specific conditions, RDDP is proved to converge to the global optimum of COPP in Section \ref{subsec:optimality_adrp_copp}. All proofs are provided in Appendix \ref{app:proofs}. 

\subsection{RDDP Process}\label{subsec:rddp_process}

Before introducing the RDDP process, the convexity of the backward reachable set $\mathcal{B}_k$ and the value function $V_k^*$ should be introduced as follows.
\begin{lemma}\label{lemma:convexity_reachset_valfunc}
    In COPP, for each $k$, the backward reachable set $\mathcal{B}_k$ is convex and compact, whereas the value function $V_k^*$ is convex in $\mathcal{B}_k$. The admissible control set $\mathcal{C}_k(\vx_k)\subset\R$ is convex and compact for each $\vx_k\in\mathcal{B}_k$; hence, $\mathcal{C}_k(\vx_k)$ is either a closed interval or a singleton. Specifically, if \eqref{eq:op_discretize_state_constraints} and \eqref{eq:op_discretize_mix_constraints} are linear inequalities, then $\mathcal{B}_k$ is a polytope.
\end{lemma}

Each loop of the standard DDP consists of a backward cut-generation pass and a forward trial pass, as shown in Fig. \ref{fig:ardp_one_loop}(a). For the $p$-th loop, the backward pass is to construct a tractable cut-based lower approximation $\mathcal{V}_k^p$ of the convex value function $V_k^*$ at each $s_k$. The forward pass is to generate a trial trajectory $(\bx_0^{p*}, \bu_0^{p*})$ guided by the approximate value function $\mathcal{V}_k^p$ and then refine the cut-approximation $\mathcal{V}_k^p$ with the trial trajectory $(\bx_0^{p*}, \bu_0^{p*})$. Different from the standard DDP, RDDP generates the backward reachable set $\mathcal{B}_k$ by BP \eqref{eq:recursion_backward} before the first loop to explicitly guarantee the feasibility of the forward trial trajectory. In this section, we assume that the backward reachable set $\mathcal{B}_k$ has been exactly computed for each $k$ in the backward reachability-generation pass. The details of the backward and forward passes are presented as follows.


\subsubsection{Backward Cut-Generation Pass}\label{sub2sec:backward_cut_pass}

At the beginning of the $p$-th loop, assume that a cut-generation set $\mathcal{S}_k^p=\{\vx_k^{i}\}_{i=1}^{I_k^p}\subset\mathcal{B}_k$ is available for each $k$, where $I_k^p\in\N^*$ is the sample number. Specifically, for $p=0$, these points can be sampled by a coarse grid or randomly. For $p>0$, the update of the cut-generation set will be analyzed in Section \ref{sub2sec:refinement}.

At the $p$-th loop, the backward pass aims to construct a tractable cut-based approximation $\mathcal{V}_k^p$ of the real value function $V_k^*$ at $s_k$. Let the terminal value be
\begin{equation}\label{eq:approximation_terminal_value}
    \mathcal{V}_N^p=\Phi=V_N^*.
\end{equation}
We aim to compute the one-step approximate value function $\widehat{\mathcal{V}}_k^p(\vx_k):\mathcal{B}_k\to\R$, where
\begin{equation}\label{eq:approximate_value_function_backward}
    \widehat{\mathcal{V}}_k^p(\vx_k) = \min_{u_k\in\mathcal{C}_k(\vx_k)} L_k(\vx_k, u_k) + \mathcal{V}_{k+1}^p(\vA_k\vx_k+\vb_k u_k).
\end{equation}
However, $\widehat{\mathcal{V}}_k^p$ remains intractable since there are uncountably many state points $\vx_k$ in $\mathcal{B}_k$. Leveraging the convexity of $V_k^*$, Section \ref{subsec:optimality_adrp_copp} will prove the convexity of $\widehat{\mathcal{V}}_k^p$. A lower approximation $\mathcal{V}_k^p$ of $\widehat{\mathcal{V}}_k^p$ is defined by optimality cuts, i.e., supporting hyperplanes of $\widehat{\mathcal{V}}_k^p$, at the cut-generation set $\mathcal{S}_k^p$. Specifically, $\forall \vx_k\in\mathcal{B}_k$, let
\begin{equation}\label{eq:support_plane}
    \mathcal{V}_k^p(\vx_k) = \max_{i=1,\ldots,I_k^p} \Big\{\widehat{\mathcal{V}}_k^p(\vx_k^i) + \vg_k^{i\top}(\vx_k-\vx_k^i)\Big\},
\end{equation}
where the cut slope $\vg_k^i\in\partial^\circ\widehat{\mathcal{V}}_k^p(\vx_k^i)$ is chosen as a Clarke subgradient \cite{clarke1990optimization} of $\widehat{\mathcal{V}}_k^p$ at $\vx_k^i$. Note that $\mathcal{C}_k(\vx_k^i)$ is an interval or singleton for each sample $\vx_k^i$. Since $\mathcal{V}_{k+1}^p$ is convex, $\widehat{\mathcal{V}}_k^p(\vx_k^i)$ can be efficiently computed through one-dimensional convex optimization \eqref{eq:approximate_value_function_backward} for each sample $\vx_k^i$.

As summarized in Fig. \ref{fig:ardp_one_loop}(b) and proved in Section \ref{subsec:optimality_adrp_copp}, it holds pointwise in $\mathcal{B}_k$ that $\mathcal{V}_k^p\leq \widehat{\mathcal{V}}_k^p\leq V_k^*$. Therefore, the tractable cut-approximation $\mathcal{V}_k^p$ also serves as a lower bound of the real value function $V_k^*$.

\begin{remark}
    In standard DDP, $\vg_k^i$ can be computed from dual multipliers of the stage optimization problem \eqref{eq:approximate_value_function_backward}, which can equivalently be interpreted as subgradients of $\widehat{\mathcal{V}}_k^p$. Instead of explicitly forming the dual of problem \eqref{eq:approximate_value_function_backward}, this paper efficiently computes the sensitivity vector $\vg_k^i$ using the Clarke subgradient. Furthermore, $\vg_k^i$ is required to be a Clarke subgradient rather than a general subgradient to ensure a tight approximation of $\widehat{\mathcal{V}}_k^p$. Consider a toy example with the state $x\in\R$. Assume that $\widehat{\mathcal{V}}_k^p(x)=x$ for $x\in\mathcal{B}_k=[0,1]$. Then, the Clarke subgradient at $x=0$ is $\partial^\circ\widehat{\mathcal{V}}_k^p(0)=\{1\}$, while the general subgradient set is $\partial\widehat{\mathcal{V}}_k^p(0)=(-\infty,1]$. In this example, we need to tightly approximate $\widehat{\mathcal{V}}_k^p$ by the optimality cut $\mathcal{V}_k^p(x)=x$ at $x=0$, while any other optimality cut, such as $\mathcal{V}_k^p(x)=0$, is loose.

    As will be shown in Section \ref{subsec:optimality_adrp_copp} and Appendix \ref{app:proofs}, the selection of Clarke subgradients is essential for the convergence and optimality guarantee of RDDP under the OPP-specific conditions on backward reachable sets.
\end{remark}

\subsubsection{Forward Trial Pass}\label{sub2sec:refinement}

In the forward trial pass for the $p$-th loop, the trial trajectory $\bx_0^{p*}$ is generated under the guidance of the approximate value function $\mathcal{V}_k^p$, as shown in Fig. \ref{fig:ardp_one_loop}(a). Assume that $\mathcal{V}_k^p$ has been constructed in the backward pass. The trial trajectory is first generated as follows:
\begin{subequations}\label{eq:forward_refine}
    \begin{align}
        &\vx_0^{p*}\in \argmin_{\vx\in\mathcal{B}_0\cap \mathcal{X}_0} \mathcal{V}_0^p(\vx),\label{eq:refine_x0}\\
        &u_k^{p*} \in \argmin_{u\in\mathcal{C}_k(\vx_k^{p*})} L_k(\vx_k^{p*}, u) + \mathcal{V}_{k+1}^p(\vA_k\vx_k^{p*}+\vb_k u),\label{eq:refine_uk}\\
        &\vx_{k+1}^{p*} = \vA_k\vx_k^{p*}+\vb_k u_k^{p*}.
    \end{align}
\end{subequations}
\eqref{eq:approximate_value_function_backward} and \eqref{eq:forward_refine} imply that for each $k$,
\begin{equation}\label{eq:widehat_V_L_V_refine}
    \widehat{\mathcal{V}}_k^p(\vx_k^{p*}) = L_k(\vx_k^{p*}, u_k^{p*}) + \mathcal{V}_{k+1}^p(\vx_{k+1}^{p*}).
\end{equation}
Then, the trial trajectory $\bx_0^{p*}$ is added to the cut-generation set for the next loop, i.e.,
\begin{equation}\label{eq:ardp_more_sample}
    \mathcal{S}_k^{p+1}\supset\mathcal{S}_k^{p}\cup\{\vx_k^{p*}\}.
\end{equation}
Note that one can also add more cut-generation points around $\vx_k^{p*}$ to further refine the approximation. Based on the updated $\mathcal{S}_k^{p+1}$, the backward cut-generation pass is performed to get $\mathcal{V}_k^{p+1}$, and then the forward trial pass is performed again to get $\bx_0^{(p+1)*}$. The above process is repeated until the trial trajectory converges or the maximum number of loops is reached.

\begin{remark}
    For each $k$, $u_k^{p*}$ and $\vx_{k+1}^{p*}$ are obtained by solving the one-dimensional convex optimization problem \eqref{eq:refine_uk}. According to \eqref{eq:refine_x0}, a convex optimization problem of dimension no greater than $m-1$ is solved to get $\vx_0^{p*}$. In practice, the initial state $\vx_0$ is usually specified, i.e., $\mathcal{X}_0=\{\vx_0\}$ is fixed. Due to the feasibility of COPP \eqref{eq:op_discretize} in Problem \ref{problem:convex}, $\vx_0\in\mathcal{B}_0$ holds. In this case, the initial state is directly set as $\vx_0^{p*}=\vx_0$ without solving any optimization problem.
\end{remark}

\begin{algorithm}[!t]
    \caption{RDDP for COPP.}
    \label{alg:ardp_copp}
    \begin{algorithmic}[1]
        \REQUIRE COPP \eqref{eq:op_discretize}, user-defined sampling scheme of  cut-generation sets, and convergence criterion;
        \ENSURE Globally optimal trajectory $\{\vx_k^*\}_{k=0}^{N}, \{u_k^*\}_{k=0}^{N-1}$;
        \STATE Initialize $p\leftarrow 0$;
        \STATE Initialize the terminal backward reachable set $\mathcal{B}_N=\mathcal{X}_\text{f}$;
        \REPEAT
            \FOR{$k\leftarrow N-1, N-2, \ldots, 0$}
                \IF{$p=0$}
                    \STATE Compute the backward reachable set $\mathcal{B}_k$ based on $\mathcal{B}_{k+1}$ via BP \eqref{eq:recursion_backward};
                    \STATE Sample $\mathcal{S}_k^0=\{\vx_k^i\}_{i=1}^{I_k^p}\subset\mathcal{B}_k$ with $I_k^p\geq1$;
                \ELSE
                    \STATE Update $\mathcal{S}_k^p\supset\mathcal{S}_k^{p-1}\cup\Big\{\vx_k^{(p-1)*}\Big\}$;
                \ENDIF
                \STATE Compute the cut-based approximated value $\widehat{\mathcal{V}}_k^p(\vx_k^i)$ and a Clarke subgradient $\vg_k^i$ at each sample $\vx_k^i\in\mathcal{S}_k^p$ based on $\mathcal{V}_{k+1}^p$ via BP \eqref{eq:approximate_value_function_backward};
                \STATE Approximate the value function based on optimality cuts at $\mathcal{S}_k^p$ and get $\mathcal{V}_k^p:\mathcal{B}_k\to\R$ via \eqref{eq:support_plane};
            \ENDFOR
            \STATE Find $\vx_0\in\mathcal{B}_0\cap \mathcal{X}_0$ to minimize $\mathcal{V}_0^p$ via \eqref{eq:refine_x0};
            \FOR{$k\leftarrow 0, 1, \ldots, N-1$}
                \STATE Find the trial control $u_k^{p*}$ to minimize $L_k(\vx_k^{p*}, u) + \mathcal{V}_{k+1}^p(\vA_k\vx_k^{p*}+\vb_k u)$ via \eqref{eq:refine_uk};
                \STATE Get $\vx_{k+1}^{p*} = \vA_k\vx_k^{p*}+\vb_k u_k^{p*}$;
            \ENDFOR
            \STATE Compute the real cost of $\{\vx_k^{p*}\}_{k=0}^{N}$ based on the real objective function \eqref{eq:op_discretize_objective}, and update the historical best trial trajectory;
        \UNTIL{the convergence criterion is satisfied;}
        \STATE \textbf{return} the historical best trial trajectory $\{\vx_k^{*}\}_{k=0}^{N}$, $\{u_k^{*}\}_{k=0}^{N-1}$ with the minimum real cost.
    \end{algorithmic}
\end{algorithm}

The proposed RDDP is summarized in Algorithm \ref{alg:ardp_copp}, where significant algorithmic flexibility is offered to allow users' tailored configurations. First, the user can compute the bidirectional reachable set and sample state points within it, which can improve the sample efficiency and achieve higher approximation accuracy. This strategy will be applied in COPP2 in Section \ref{subsec:ardp_copp2}. Second, the user can choose different sampling schemes for each loop. The only requirement is that the trial trajectory of the previous loop, i.e., $\vx_k^{(p-1)*}$, should be included in the cut-generation set $\mathcal{S}_k^p$ for each $p,k$. Third, the approximation value $\mathcal{V}_k^p(\vx_k^i)$ and its Clarke subgradient might keep unchanged across different loops, which can be reused to save computational costs. Finally, the convergence criterion is recommended to define based on the DDP optimality gap which will be introduced in Theorem \ref{thm:near_optimality_gap}.

\subsection{Global Optimality Guarantee}\label{subsec:optimality_adrp_copp}

This section proves the convergence and optimality of RDDP for COPP under the OPP-specific condition. Before that, we first investigate the existence of Clarke subgradient and the Lipschitz continuity of $\widehat{\mathcal{V}}_k^p$ in \eqref{eq:approximate_value_function_backward} for well-posedness. Different from the standard DDP under the relatively complete recourse assumption, this paper only requires the following regularity condition on the admissible control set.

\begin{assumption}\label{assum:Lipschitz}
    For each $k$, the admissible control set $\mathcal{C}_k(\vx_k)$ can be represented as 
    \begin{equation}
    \mathcal{C}_k(\vx_k) = \{u\in\R \mid \underbar{u}_k(\vx_k)\leq u\leq \bar{u}_k(\vx_k)\},
    \end{equation}
    where $\underbar{u}_k$ and $-\bar{u}_k$ are convex. Assume that $\underbar{u}_k$ and $\bar{u}_k$ are both locally Lipschitz continuous in their domains.
\end{assumption}

\begin{remark}
    Assumption \ref{assum:Lipschitz} is typically satisfied in OPP. In the widely studied OPP2 and OPP3, the kinodynamic constraints and their convexified variants are often affine in state and control \cite[Appendix A.2]{wang2026online}. In this case, $\underbar{u}_k$ and $\bar{u}_k$ are piecewise linear and thus locally Lipschitz continuous. Only in degenerate cases can the local Lipschitz continuity fail in the boundary of $\mathcal{B}_k$. For example, if $u_k\leq\bar{u}_k(\vx_k)=\sqrt{1-x_{k,1}^2}$, then $\bar{u}_k$ is not locally Lipschitz continuous at boundary points $x_{k,1}=\pm1$. For such an uncommon case, one can consider a slightly contracted feasible domain to exclude the singularity. Therefore, Assumption \ref{assum:Lipschitz} is not restrictive for the OPP instances considered in this paper.
\end{remark}

\begin{proposition}\label{prop:Lipschitz}
    For each $k$, $\widehat{\mathcal{V}}_k^p$ is convex and locally Lipschitz continuous in $\mathcal{B}_k$ with the local Lipschitz constant independent of $p$. In other words, $\forall\vx_k\in\mathcal{B}_k$, $\exists \varepsilon_k^V(\vx_k)>0$ small enough and $C_k^V(\vx_k)>0$ independent of $p$, s.t. $\forall p\geq0$, $\vy_1,\vy_2\in\mathcal{B}_k\cap\mathbb{B}(\vx_k,\varepsilon_k^V(\vx_k))$, we have
    \begin{equation}
        \abs{\widehat{\mathcal{V}}_k^p(\vy_1)-\widehat{\mathcal{V}}_k^p(\vy_2)}\leq C_k^V(\vx_k)\norm[2]{\vy_1-\vy_2},
    \end{equation}
    where $\mathbb{B}(\vx_k,\varepsilon_k^V(\vx_k))$ is the Euclidean ball centered at $\vx_k$ with radius $\varepsilon_k^V(\vx_k)$.  Furthermore, the Clarke subgradient set $\partial^\circ\widehat{\mathcal{V}}_k^p(\vx_k)$ is non-empty and bounded by $C_k^V(\vx_k)$. In other words, $\forall \vg_k\in\partial^\circ\widehat{\mathcal{V}}_k^p(\vx_k)$, we have $\norm[2]{\vg_k}\leq C_k^V(\vx_k)$.

    The above conclusion also holds for $\mathcal{V}_k^p$ and the real value function $V_k^*$ with the same local Lipschitz constant $C_k^V(\vx_k)$.
\end{proposition}

Proposition \ref{prop:Lipschitz} guarantees the well-posedness of the RDDP process under OPP-specific settings. In the following analysis, Proposition \ref{prop:bound_of_value_function}, Theorems \ref{prop:bound_of_value_function} and \ref{thm:near_optimality_gap} recover a DDP-like upper-lower-bound convergence structure under OPP-specific feasibility and regularity conditions. The assumptions of RDDP, however, are different from those in classical DDP. Classical DDP usually establishes such results under relatively complete recourse over a prescribed state domain. In contrast, RDDP explicitly uses the backward reachable sets $\mathcal{B}_k$ as the domain of the value functions and enforces recourse feasibility through the constraint $\vx_{k+1}\in\mathcal{B}_{k+1}$. As a result, the optimality theory below is built on OPP-specific reachable-domain conditions and local Lipschitz regularity of the admissible-control boundaries, rather than on a global recourse assumption over the entire state space.

\begin{proposition}\label{prop:bound_of_value_function}
    In the backward cut-generation pass of RDDP, for each $k$ and $p$, consider any feasible solution $\bx_k=(\vx_j)_{j=k}^N$, $\bu_k=(u_j)_{j=k}^{N-1}$ of COPP \eqref{eq:op_discretize}. We have
    \begin{equation}\label{eq:bound_of_value_function}
        \mathcal{V}_k^p(\vx_k) \leq\widehat{\mathcal{V}}_k^p(\vx_k) \leq V_k^*(\vx_k) \leq J_k(\bx_k,\bu_k),
    \end{equation}
    where $J_k(\bx_k,\bu_k)$ is the real cost in \eqref{eq:op_discretize_objective}, i.e.,
    \begin{equation}
        J_k(\bx_k,\bu_k) = \sum_{j=k}^{N-1} L_j(\vx_j, u_j) + \Phi(\vx_N).
    \end{equation}
\end{proposition}

In \eqref{eq:bound_of_value_function}, the piecewise linear approximation $\mathcal{V}_k^p$ and the real cost $J_k$ are both computationally tractable, whereas the real value function $V_k^*$ and the one-step approximation $\widehat{\mathcal{V}}_k^p$ are intractable. As illustrated in Fig. \ref{fig:ardp_one_loop}(b), the key theoretical mechanism is a tractable upper-lower bound certificate: the real cost $J_k$ provides an upper bound on the optimal value, while the cut-based approximate $\mathcal{V}_k^p$ provides a lower bound.

\begin{theorem}\label{thm:near_optimality_gap}
    For the $p$-th loop of RDDP, consider the feasible trial trajectory $\bx_0^{p*}=(\vx_k^{p*})_{k=0}^N$ and $\bu_0^{p*}=(u_k^{p*})_{k=0}^{N-1}$ generated in the forward pass. Denote the globally optimal value of COPP \eqref{eq:op_discretize} by $J^*$. Then, the optimality gap is bounded by
    \begin{equation}
        0\leq J_0(\bx_0^{p*},\bu_0^{p*}) - J^* \leq \varepsilon_\text{gap}(\bx_0^{p*},\bu_0^{p*}),
    \end{equation}
    where the \textit{DDP optimality gap} $\varepsilon_\text{gap}$ is defined as
    \begin{equation}\label{eq:dual_gap}
        \varepsilon_\text{gap}(\bx_0^{p*},\bu_0^{p*})=J_0(\bx_0^{p*},\bu_0^{p*}) - \mathcal{V}_0^p(\vx_0^{p*}).
    \end{equation}
\end{theorem}

Theorem \ref{thm:near_optimality_gap} provides a stopping criterion for Algorithm \ref{alg:ardp_copp}, i.e., the DDP optimality gap $\varepsilon_\text{gap}$ is smaller than a user-defined threshold. Finally, the global optimality convergence of RDDP is presented as follows.

\begin{theorem}\label{thm:adrp_optimal}
    The historical best cost of RDDP converges to the globally optimal value. Specifically,
    \begin{equation}\label{eq:convergence_condition}
        \liminf_{p \to \infty} \varepsilon_\text{gap}(\bx_0^{p*},\bu_0^{p*}) = 0,
    \end{equation}
    where notations in Theorem \ref{thm:near_optimality_gap} are applied. Furthermore, each accumulation point of the trial sequence $\{(\bx_0^{p*},\bu_0^{p*})\}_{p=0}^\infty$ generated by RDDP is a globally optimal solution of COPP \eqref{eq:op_discretize}. In other words, any convergent subsequence, the existence of which is guaranteed by the compactness of the feasible set, converges to a globally optimal solution of COPP \eqref{eq:op_discretize}.
\end{theorem}

Based on Theorem \ref{thm:adrp_optimal}, the proposed RDDP can be regarded as a principled alternative to conventional CO methods for COPP. Note that RDDP attains the same global optimality while replacing the large-scale optimization problem \eqref{eq:op_discretize} by a sequence of low-dimensional ones, thereby enabling substantially more efficient computation, as verified in Section \ref{subsec:numerical_experiments}.

\begin{remark}
    Proposition \ref{prop:bound_of_value_function} and Theorem \ref{thm:adrp_optimal} provide a key insight which enables adaptive refinement in RDDP. Specifically, we define a tractable one-step DDP optimality gap as follows:
    \begin{equation}\label{eq:dual_gap_k}
        \varepsilon_\text{gap}^k(\vx_k^{p*}, u_k^{p*}) = L_k(\vx_k^{p*}, u_k^{p*}) + \mathcal{V}_{k+1}^p(\vx_{k+1}^{p*}) - \mathcal{V}_k^p(\vx_k^{p*}).
    \end{equation}
    If $\varepsilon_\text{gap}^k(\vx_k^{p*}, u_k^{p*})$ is large, then the current cut-based lower approximation $\mathcal{V}_k^p$ is loose around $\vx_k^{p*}$. We can add more cut-generation points around $\vx_k^{p*}$ to refine the approximation for the $(p+1)$-th loop, which can accelerate convergence.
\end{remark}

\begin{remark}
    The global optimality of RDDP is guaranteed by an assumption that the backward reachable set $\mathcal{B}_k$ is exactly computed by \eqref{eq:recursion_backward} for each $k$. As will be introduced in Section \ref{subsec:ardp_copp3}, only conservatively approximate $\mathcal{B}_k$ is available for COPP3 \cite{dio2025time}. As will be proved in Section \ref{subsec:KKT_GOPP_Conservative}, the global optimality of the original COPP problem \eqref{eq:op_discretize} can also be attained by iteratively updating reachable sets when conservative inner approximations of reachable sets are used.
\end{remark}

\section{RDDP for GOPP with KKT Convergence}\label{sec:RDDP_GOPP}

In Section \ref{sec:RDDP_COPP}, we have proposed RDDP, where global optimality of COPP is guaranteed by Theorem \ref{thm:adrp_optimal}. In practice, some objectives and constraints in OPP \eqref{eq:op_discretize} are non-convex, such as the jerk constraint \eqref{eq:jerk_constraint} in OPP3. However, the standard DDP is primarily designed for convex multistage problems, whereas cut-based approximations are not guaranteed to be global lower bounds of value functions without convexity.

Under OPP-specific conditions, this section extends RDDP to GOPP, i.e., the non-convex Problem \ref{problem:general}, through embedding RDDP into the SCP framework. Since the theoretical results in Section \ref{subsec:optimality_adrp_copp} are established under the assumption that backward reachable sets are exactly computed, we first analyze the KKT convergence of RDDP for GOPP under such assumption in Section \ref{subsec:KKT_GOPP_exact_backward}. In Section \ref{subsec:KKT_GOPP_Conservative}, we extend the results to the case where only conservative approximation of the backward reachable set is available, which is highly common in OPP3. All proofs are given in Appendix \ref{app:proofs}.

\subsection{KKT Convergence with Exact Reachable Sets}\label{subsec:KKT_GOPP_exact_backward}

Leveraging the global optimality of RDDP in COPP, this section directly applies DC decomposition to the non-convex functions in GOPP. For each iteration $q$ of DC decomposition, we can construct a convex optimization problem by convexifying the non-convex functions, which can be solved by RDDP. In this section, we assume that the backward reachable sets of each convexified problem are exactly computed by \eqref{eq:recursion_backward}.

According to assumptions in Problem \ref{problem:general}, all functions $L_k$, $\Phi$, $\vf_k$, $\vh_k$, and boundary sets $\mathcal{X}_0$, $\mathcal{X}_\text{f}$ in GOPP \eqref{eq:op_discretize} are DC-decomposable. For the $q$-th iteration of DC decomposition, $q\geq0$, assume that the reference point $\bx_0^{(q)}=\Big(\vx_k^{(q)}\Big)_{k=0}^{N}$, $\bu_0^{(q)}=\Big(u_k^{(q)}\Big)_{k=0}^{N-1}$ is given. Then, we construct the following convex optimization subproblem by convexifying the non-convex functions and sets at $\bx_0^{(q)},\bu_0^{(q)}$:
\begin{subequations}\label{eq:op_discretize_convexity}
    \begin{align}
        \min_{\vx_k,u_k} \quad &\sum_{k=0}^{N-1} \overline{L}_k\Big(\vx_k, u_k; \vx_k^{(q)},u_k^{(q)}\Big) + \overline{\Phi}\Big(\vx_N;\vx_N^{(q)}\Big) \label{eq:op_discretize_convexity_objective}\\
        \st \quad & \vx_{k+1} = \vA_k\vx_k + \vb_k u_k,\label{eq:op_discretize_convexity_dynamics} \\
        &\overline{\vf}_k\Big(\vx_k;\vx_k^{(q)}\Big) \leq \vzero, \label{eq:op_discretize_convexity_state_constraints}\\
        &\overline{\vh}_k\Big(\vx_k, u_k; \vx_k^{(q)}, u_k^{(q)}\Big) \leq \vzero,\label{eq:op_discretize_convexity_mix_constraints}\\
        &\vx_0\in\overline{\mathcal{X}}_0\Big(\vx_0^{(q)}\Big), \,\,\vx_N\in\overline{\mathcal{X}}_\text{f}\Big(\vx_N^{(q)}\Big),\label{eq:op_discretize_convexity_boundary_constraints}
    \end{align}
\end{subequations}
where the convexified functions and sets are defined in Definition \ref{def:DC_decomposable}. Note that the convexified feasible set of \eqref{eq:op_discretize_convexity} is a subset of the original feasible set of GOPP \eqref{eq:op_discretize}, as shown in Fig. \ref{fig:ardp_gopp}. Apply RDDP to solve the convex subproblem \eqref{eq:op_discretize_convexity} and obtain an optimal solution $\bx_0^{(q+1)}=\Big(\vx_k^{(q+1)}\Big)_{k=0}^{N}$, $\bu_0^{(q+1)}=\Big(u_k^{(q+1)}\Big)_{k=0}^{N-1}$ which is used as the reference point for the next iteration of DC decomposition. Then, the following theorem guarantees the KKT convergence of the above DC decomposition process.

\begin{figure}[!t]
    \centering
    \includegraphics[width=0.85\columnwidth]{ardp_gopp.png}
    \caption{Relation between the backward reachable sets of the convexified problem and the original GOPP.}
    \label{fig:ardp_gopp}
\end{figure}

\begin{theorem}\label{thm:KKT_exact_reachable_set}
    Consider GOPP \eqref{eq:op_discretize} and the above DC decomposition process. An initial reference point $\bx_0^{(0)}=\Big(\vx_k^{(0)}\Big)_{k=0}^{N}$, $\bu_0^{(0)}=\Big(u_k^{(0)}\Big)_{k=0}^{N-1}$ is given. Assume that the convexified problem \eqref{eq:op_discretize_convexity} for $q=0$ is feasible. Then, $\forall q\geq0$, the convexified problem \eqref{eq:op_discretize_convexity} remains feasible. The real cost of the solution sequence $\Big\{\bx_0^{(q)},\bu_0^{(q)}\Big\}_{q=0}^\infty$ is non-increasing, i.e.,
    \begin{equation}
        J_0\Big(\bx_0^{(q+1)},\bu_0^{(q+1)}\Big) \leq J_0\Big(\bx_0^{(q)},\bu_0^{(q)}\Big),
    \end{equation}
    where $J_0\Big(\bx_0^{(q)},\bu_0^{(q)}\Big)=\sum_{k=0}^{N-1} L_k\Big(\vx_k^{(q)}, u_k^{(q)}\Big)+\Phi\Big(\vx_N^{(q)}\Big)$.

    Furthermore, assume that $\forall q\geq0$, the convexified problem \eqref{eq:op_discretize_convexity} satisfies certain constraint qualification (CQ) \cite{boyd2004convex}, such as the Slater's condition and linear independence constraint qualification (LICQ). Then, each accumulation point $(\bx_0^*,\bu_0^*)$ of the solution sequence $\Big\{\Big(\bx_0^{(q)},\bu_0^{(q)}\Big)\Big\}_{q=0}^\infty$ is a KKT solution of the original GOPP \eqref{eq:op_discretize} if $(\bx_0^*,\bu_0^*)$ also satisfies the CQ condition of GOPP \eqref{eq:op_discretize}.
\end{theorem}

\begin{algorithm}[!t]
    \caption{RDDP for GOPP.}
    \label{alg:ardp_gopp}
    \begin{algorithmic}[1]
        \REQUIRE GOPP \eqref{eq:op_discretize}, convergence criterion of DC decomposition, initial reference point $\Big(\bx_0^{(0)}, \bu_0^{(0)}\Big)$, and inputs of Algorithm \ref{alg:ardp_copp};
        \ENSURE KKT trajectory $\{\vx_k^*\}_{k=0}^{N}, \{u_k^*\}_{k=0}^{N-1}$;
        \STATE Initialize $q\leftarrow 0$;
        \REPEAT
            \STATE Convexify GOPP \eqref{eq:op_discretize} at $\bx_0^{(q)},\bu_0^{(q)}$ to construct the convex optimization problem \eqref{eq:op_discretize_convexity} or \eqref{eq:op_discretize_convexity_conservative}, where the computed backward reachable sets $\mathcal{B}_{\text{c}k}^{(q)}$ should satisfy the condition \eqref{eq:KKT_approximate_reachable_set_condition} \textbf{if} $q\geq1$;
            \STATE Apply Algorithm \ref{alg:ardp_copp} to solve problems \eqref{eq:op_discretize_convexity} or \eqref{eq:op_discretize_convexity_conservative}, and obtain the optimal solution $\Big(\bx_0^{(q+1)}, \bu_0^{(q+1)}\Big)$;
            \STATE Update $q\leftarrow q+1$;
        \UNTIL{the convergence criterion is satisfied;}
        \STATE \textbf{return} the last solution $\Big(\bx_0^{(q)}, \bu_0^{(q)}\Big)$.
    \end{algorithmic}
\end{algorithm}

\subsection{KKT Convergence with Conservative Reachable Sets}\label{subsec:KKT_GOPP_Conservative}

In practice, exact backward reachable sets of each convexified problem may be unavailable. For example, backward reachable sets of COPP3 are typically approximated by a polytope with limited number of edges \cite{dio2025time}. This section proves that the KKT convergence of the DC decomposition process is still guaranteed with conservative reachable sets. The approximate backward reachable set can be regarded as some additional state constraints $\vr^{(q)}_k(\vx_k)\leq \vzero$ in the convexified subproblem \eqref{eq:op_discretize_convexity} for each iteration $q$. The considered convexified subproblem for the $q$-th iteration is as follows:
\begin{subequations}\label{eq:op_discretize_convexity_conservative}
    \begin{align}
        \min_{\vx_k,u_k} \quad &\sum_{k=0}^{N-1} \overline{L}_k\Big(\vx_k, u_k; \vx_k^{(q)},u_k^{(q)}\Big) + \overline{\Phi}\Big(\vx_N;\vx_N^{(q)}\Big)\label{eq:op_discretize_convexity_conservative_objective}\\
        \st \quad & \vr^{(q)}_k(\vx_k)\leq \vzero,\,\,k=0,\ldots,N-1,\label{eq:op_discretize_convexity_conservative_additional_constraints}\\
        & \text{Constraints \eqref{eq:op_discretize_convexity_dynamics}--\eqref{eq:op_discretize_convexity_boundary_constraints}}.\label{eq:op_discretize_convexity_conservative_other_constraints}
    \end{align}
\end{subequations}
Then, an optimal solution $\Big(\bx_0^{(q+1)},\bu_0^{(q+1)}\Big)$ of the conservative problem \eqref{eq:op_discretize_convexity_conservative} can be obtained by RDDP. The following theorem guarantees the KKT convergence of RDDP, where only an additional local exactness condition is required that the extra state constraints $\vr^{(q)}_k(\vx_k)\leq \vzero$ should not influence the feasibility of the convexified problem in a neighborhood of the reference point $\Big(\bx_0^{(q)},\bu_0^{(q)}\Big)$, as shown in Fig. \ref{fig:ardp_gopp}. The RDDP process of GOPP is summarized in Algorithm \ref{alg:ardp_gopp}.

\begin{theorem}\label{thm:KKT_approximate_reachable_set}
    Consider GOPP \eqref{eq:op_discretize} and the above DC decomposition process. An initial reference point $\Big(\bx_0^{(0)},\bu_0^{(0)}\Big)$ is given. For the $q$-th iteration, denote the exact and conservative backward reachable sets as $\mathcal{B}_{\text{e}k}^{(q)}$ and $\mathcal{B}_{\text{c}k}^{(q)}$, respectively. In other words, $\mathcal{B}_{\text{e}k}^{(q)}$ is the intractable theoretical backward reachable set determined by constraints \eqref{eq:op_discretize_convexity_conservative_other_constraints}, whereas $\mathcal{B}_{\text{c}k}^{(q)}\subset\mathcal{B}_{\text{e}k}^{(q)}$ is the tractable backward reachable set determined by constraints \eqref{eq:op_discretize_convexity_conservative_additional_constraints} and \eqref{eq:op_discretize_convexity_conservative_other_constraints}. Assume that $\exists \delta>0$ independent of $q$, s.t.
    \begin{equation}\label{eq:KKT_approximate_reachable_set_condition}
        \forall q\geq1,\,\mathcal{B}_{\text{c}k}^{(q)}\cap\mathbb{B}\Big(\vx_k^{(q)},\delta\Big)=\mathcal{B}_{\text{e}k}^{(q)}\cap\mathbb{B}\Big(\vx_k^{(q)},\delta\Big).
    \end{equation}
    Then, the feasibility, non-increasing real cost of the solution sequence, and the KKT convergence of the RDDP process are guaranteed under the same conditions as in Theorem \ref{thm:KKT_exact_reachable_set}.
\end{theorem}

Theorem \ref{thm:KKT_approximate_reachable_set} can be reduced to COPP3 as follows. Specifically, the DC decomposition process is reduced to a local update of the approximate reachable sets around the reference point, while all functions and sets are not convexified.

\begin{corollary}\label{cor:global_optimal_approximate_reachable_set}
    Consider COPP where only conservative approximation of the backward reachable sets is available. Following the same conditions and process in Theorem \ref{thm:KKT_approximate_reachable_set}, RDDP converges to a globally optimal solution.
\end{corollary}

\section{RDDP for OPP2 and OPP3}\label{sec:ardp_opp2_opp3}

This section develops efficient instantiations of RDDP for OPP2 and OPP3. Specifically, this section focuses on COPP since GOPP can be solved by the process in Section \ref{sec:RDDP_GOPP}.

\subsection{RDDP for COPP2}\label{subsec:ardp_copp2}

For COPP2, the backward reachable sets can be computed exactly as intervals by existing RA techniques. The one-dimensional state is denoted by $\vx_k=x_k\triangleq x_{k,1}=\frac12\dot{s}^2(s_k)\in\R_+$. Then, the dynamic equation \eqref{eq:op_discretize_dynamics} implies that $x_{k+1}=x_k+u_k\Delta_k$.

\subsubsection{Bidirectional Reachable Set}\label{sub2sec:ardp_opp2_reachable_set}
We can directly apply the RA approach in \cite{pham2018new} for COPP2. To improve the sampling efficiency of RDDP, we choose to compute the bidirectional reachable set instead of the backward reachable set. Note that such modification does not influence the optimality and convergence of RDDP, as the bidirectional reachable set is the subset of the backward reachable set of those states that can be reached from the initial state set $\mathcal{X}_0$. The backward reachable set $\mathcal{B}_{k}'=[\underbar{x}_{k}',\bar{x}_{k}']$ is computed via BP in \eqref{eq:recursion_backward}, as shown in Fig. \ref{fig:reachable_set_copp2}(a). Then, as shown in Fig. \ref{fig:reachable_set_copp2}(b), the bidirectional reachable set $\mathcal{B}_k=[\underbar{x}_k,\bar{x}_k]$ is computed via FP by solving the following two-dimensional convex optimization problem:
\begin{subequations}
    \begin{align}
        \minmax_{x_k,u_{k-1}}\quad&x_k\\
        \st\quad&\vh_{k-1}(x_k-u_{k-1}\Delta_{k-1},u_{k-1})\leq\vzero,\\
        &\underbar{x}_{k-1}\leq x_k-u_{k-1}\Delta_{k-1}\leq \bar{x}_{k-1},\\
        &\underbar{x}_k'\leq x_k\leq \bar{x}_k',
    \end{align}
\end{subequations}
where $\mathcal{B}_{0}=\mathcal{B}_{0}'\cap\mathcal{X}_0$.

\begin{figure}[!]
    \centering
    \includegraphics[width=0.87\columnwidth]{reachable_set_copp2.png}
    \caption{Computation of reachable sets for COPP2. (a) Backward propagation of the backward reachable set $\mathcal{B}_k'$. (b) Forward propagation of the bidirectional reachable set $\mathcal{B}_k$.}
    \label{fig:reachable_set_copp2}
\end{figure}

\begin{figure*}[!]
    \centering
    \includegraphics[width=\textwidth]{reachable_set_copp3_extend.png}
    \caption{Computation of backward reachable sets for COPP3 based on projection. (a) The one-step exact reachable set $\mathcal{B}_{\text{e}k}$ with sampled vertices. (b) The proposed extension-based method in Section \ref{sub2sec:backward_set_copp3_extend}. (c) The baseline in \cite{dio2025time}. The proposed method provides a more accurate estimation of the residual area.}
    \label{fig:reachable_set_copp3_extend}
\end{figure*}

\subsubsection{Backward Cut-Generation Pass}\label{sub2sec:ardp_opp2_subgradient}

The backward cut-generation pass of RDDP for COPP2 follows the standard procedure in Section \ref{sub2sec:backward_cut_pass}. The admissible control set $\mathcal{C}_k(x_k)=[\underbar{u}_k(x_k),\bar{u}_k(x_k)]$ can be solved by the one-dimensional convex optimization problem \eqref{eq:admissable_control_set}, where the active constraints can be used to compute the directional derivatives $\frac{\mathrm{d}\underbar{u}_k}{\mathrm{d}x_k}(x_k)$ and $\frac{\mathrm{d}\bar{u}_k}{\mathrm{d}x_k}(x_k)$. Then, the cut-based approximation $\mathcal{V}_k^p$ is constructed by the one-dimensional convex problem \eqref{eq:approximate_value_function_backward}.

\subsubsection{Forward Trial Pass} In the forward pass, approximation value functions $\mathcal{V}_k^p$ have been computed for all $k$. Following the standard process in Section \ref{sub2sec:refinement}, the trial trajectory is generated by solving a sequence of one-dimensional convex optimization problems \eqref{eq:forward_refine}.

Since the reachable set $\mathcal{B}_k$ is exactly computed for COPP2, the global optimality convergence of the trial trajectories with respect to (w.r.t.) the original COPP \eqref{eq:op_discretize} is guaranteed by Theorem \ref{thm:adrp_optimal}.

\subsection{RDDP for COPP3}\label{subsec:ardp_copp3}

In this section, COPP3 typically denotes the convexified subproblem of GOPP3 with non-convex 3rd-order constraints. The backward cut-generation pass and forward trial pass follow the standard process in Section \ref{subsec:rddp_process}, where a sequence of one- and two-dimensional convex optimization problems are solved. Therefore, this section focuses on the computation of backward reachable sets for COPP3. In fact, a challenge for RA-based COPP3 approaches is the computation and representation of the backward reachable sets, which are generally convex but intractable sets in $\R^2$. According to \eqref{eq:backward_reachable_definition}, the exact $\mathcal{B}_k$ is the projection of a high-dimensional convex set represented by numerous convex constraints. Specifically, as the BP process goes backward and $k$ decreases, the number of constraints for representing $\mathcal{B}_k$ grows. As a result, the number of active constraints at the boundary of $\mathcal{B}_k$ grows uncontrollably, which leads to the intractability of $\mathcal{B}_k$.

Leveraging the recent work \cite{dio2025time} and our established Theorem \ref{thm:KKT_approximate_reachable_set}, we can compute the backward reachable sets by a conservative polytope with limited number of edges, while guaranteeing the global optimality for COPP3 or the KKT convergence for GOPP3 under local exactness conditions. Specifically, this section presents two methods for computing conservative backward reachable sets for efficiency under different conditions. The improved reachable-set computation can also be plugged into TOPP3-RA solvers, leading to the TOPP3-RA-IRC variant evaluated in this paper.

In contrast to COPP2, we recommend to use backward reachable sets instead of bidirectional reachable sets in COPP3, as the approximation of bidirectional reachable sets would significantly increase the conservativeness and hinder the optimality. In COPP3, the two-dimensional state is denoted by $\vx_k=(x_k,y_k)\triangleq (x_{k,1},x_{k,2})=(\frac12\dot{s}^2(s_k),\ddot{s}(s_k))\in\R_+\times\R$. Then, the dynamic equation \eqref{eq:op_discretize_dynamics} implies $x_{k+1}=x_k+y_k\Delta_k+\frac12u_k\Delta_k^2$ and $y_{k+1}=y_k+u_k\Delta_k$. In this section, we still denote the approximate backward reachable set by $\mathcal{B}_k$ for brevity.

\subsubsection{Projection-Based Backward Reachable Set}\label{sub2sec:backward_set_copp3_extend}

\begin{figure*}[!]
    \centering
    \includegraphics[width=\textwidth]{reachable_set_copp3_clip.png}
    \caption{Computation of the one-step exact backward reachable set $\mathcal{B}_{\text{e}k}$ for COPP3 based on searching. (a) The geometric projection relation between $\mathcal{B}_{\text{e}k}$ and $\widehat{\mathcal{B}}_k$. (b) Searching along a line. (c) Searching along a plane.}
    \label{fig:reachable_set_copp3_clip}
\end{figure*}

This section introduces an improved method of \cite[Section \RomanNum{3}-B2]{dio2025time}. Let $\mathcal{B}_{\text{c}N}=\mathcal{X}_\text{f}$. Given a convex set $\mathcal{B}_{\text{c}\,k+1}$, the one-step exact backward reachable set $\mathcal{B}_{\text{e}k}$ is the projection of the following three-dimensional convex set:
\begin{align}
    \widehat{\mathcal{B}}_k=&\Big\{(x_k, y_k, u_k)\mid \vf_k(x_k,y_k)\leq\vzero,\,\vh_k(x_k,y_k,u_k)\leq\vzero, \notag\\
    &\Big(x_k+y_k\Delta_k+\frac12u_k\Delta_k^2,\, y_k+u_k\Delta_k\Big)\in\mathcal{B}_{\text{c}\,k+1}\Big\}.\label{eq:reachable_set_copp3_3d_before_projection}
\end{align}
To project a boundary point of $\mathcal{B}_{\text{e}k}$ to a vertex of the conservative reachable set $\mathcal{B}_{\text{c}k}$, we can solve the following three-dimensional convex optimization problem:
\begin{equation}\label{eq:reachable_set_copp3_extend}
    (x_k^i, y_k^i, u_k^i)\in\argmax_{(x_k,y_k,u_k)\in\widehat{\mathcal{B}}_k} a_{i}x_k+ b_{i}y_k,
\end{equation}
where $(a_i,b_i)\not=\vzero$ is the sampling direction. 

As shown in Fig. \ref{fig:reachable_set_copp3_extend}(a), the exact $\mathcal{B}_{\text{e}k}$ can be obtained by sampling infinitely many directions, which is intractable. In practice, an upper bound of the number of vertices $P_k$ is given. Then, we can sample no more than $P_k$ vertices of $\mathcal{B}_{\text{e}k}$ with different sampling directions $(a_i,b_i)$. As a subset of $\mathcal{B}_{\text{e}k}$, the convex hull of the sampled vertices is computed as the conservative backward reachable set $\mathcal{B}_{\text{c}k}$.

The remaining issue is how to choose the sampling directions. Inspired by \cite{dio2025time}, the area of $\mathcal{B}_{\text{c}k}$ is heuristically maximized to minimize the conservativeness. Assume that there has been $i<P_k$ vertices sampled, denoted by $\{\vx_k^j=(x_k^j, y_k^j)\}_{j=1}^i$. As shown in Fig. \ref{fig:reachable_set_copp3_extend}(b1), denote $\mathcal{B}_\text{in}$ by the convex hull of the sampled vertices, which is an inner convex polytopic approximation of $\mathcal{B}_{\text{e}k}$. For each sampled vertex $\vx_k^j$, we can compute the left and right tangent lines of $\mathcal{B}_{\text{e}k}$ at $\vx_k^j$ by determining the active constraints of the above optimization problem. As shown in Fig. \ref{fig:reachable_set_copp3_extend}(b2), denote $\mathcal{B}_{\text{out}}$ by the region enclosed by these tangent lines, which is an outer convex polytopic approximation of $\mathcal{B}_{\text{e}k}$. Hence,
\begin{equation}
    \mathcal{B}_\text{in}\subseteq\mathcal{B}_{\text{e}k}\subseteq\mathcal{B}_\text{out}.
\end{equation}
Although $\mathcal{B}_{\text{e}k}$ is intractable, the polytopes $\mathcal{B}_\text{in}$ and $\mathcal{B}_\text{out}$ are both tractable. Denote the vertices of $\mathcal{B}_\text{out}$ by $\{\hat{\vx}_k^j\}_{j=1}^i$. Then, $\Setminus{\mathcal{B}_\text{out}}{\mathcal{B}_\text{in}}$ is the union of triangles formed by $\hat{\vx}_k^j$, $\vx_k^j$, and $\vx_k^{j+1}$, where $\vx_k^{i+1}$ refers to $\vx_k^1$. Similar to \cite{dio2025time}, we select the triangle with the largest area, i.e., the maximal residual area. Then, the next vertex is sampled with the outer normal direction of the corresponding segment $\vx_k^j$--$\vx_k^{j+1}$. The above process is repeated until $P_k$ different vertices are sampled or $\mathcal{B}_{\text{e}k}$ is exactly obtained. Finally, let $\mathcal{B}_{\text{c}k}$ be the convex hull of all sampled vertices.

\begin{remark}
    The difference between the proposed method and \cite{dio2025time} lies in the construction of the outer polytope $\mathcal{B}_\text{out}$. In \cite{dio2025time}, the outer polytope is constructed by the support lines of \eqref{eq:reachable_set_copp3_extend}, i.e., the orthogonal lines to $(a_i,b_i)$, as shown in Fig. \ref{fig:reachable_set_copp3_extend}(c). In contrast, our method constructs $\mathcal{B}_\text{out}$ by the two-sided tangent lines at each sampled vertex, which provides a more accurate estimation of the residual area. As shown in Fig. \ref{fig:reachable_set_copp3_extend}(b)-(c), the proposed method empirically achieves a tighter approximation of $\mathcal{B}_{\text{e}k}$ than \cite{dio2025time}, which contributes to a better optimality of the generated trajectory.
\end{remark}

\subsubsection{Searching-Based Backward Reachable Set}\label{sub2sec:backward_set_copp3_clip}

This section proposes a more efficient method for computing the backward reachable set for the case where all constraints in COPP3 \eqref{eq:op_discretize} are linear. Note that the linearity of constraints is a common case in practice \cite[Appendix A.2]{wang2026online}. By Lemma \ref{lemma:convexity_reachset_valfunc}, the exact backward reachable set is a convex polytope in $\R^2$. As shown in Fig. \ref{fig:reachable_set_copp3_clip}(a), the one-step exact reachable set $\mathcal{B}_{\text{e}k}$ is the projection of a three-dimensional convex polytope $\widehat{\mathcal{B}}_k$ in \eqref{eq:reachable_set_copp3_3d_before_projection}.

The key idea is to directly compute the exact $\mathcal{B}_{\text{e}k}$ by iterative vertex searching. Specifically, the first vertex $\vx_k^1$ of $\mathcal{B}_{\text{e}k}$ is obtained through the three-dimensional linear programming (LP) \eqref{eq:reachable_set_copp3_extend} with an arbitrary sampling direction. Assume that the last sampled three-dimensional vertex is $(x_k^i, y_k^i, u_k^i)$ with left and right tangent directions $\vd_{i}^-$ and $\vd_i^+$. Then, we can search for the next vertex along the right tangent direction $\vd_i^+$ which corresponds to either an edge or a plane of $\widehat{\mathcal{B}}_k$. As shown in Fig. \ref{fig:reachable_set_copp3_clip}, a one- or two-dimensional LP is solved to obtain the next vertex for the two cases, respectively.

Finally, the standard ear-clipping method is applied to the exact $\mathcal{B}_{\text{e}k}$ to reduce the number of vertices to be no more than $P_k$ for efficiency. Specifically, the ear-clipping method iteratively removes the triangle of the smallest area formed by three adjacent vertices until the number of vertices is no more than $P_k$. The resulting convex hull of the remaining vertices is the conservative backward reachable set $\mathcal{B}_{\text{c}k}$.

\begin{remark}
    Intuitively, the projection-based method in Section \ref{sub2sec:backward_set_copp3_extend} can save computational cost when $P_k$ is small, whereas the searching-based method in this section needs to fully solve all vertices of $\mathcal{B}_{\text{e}k}$. However, the searching-based method is much more efficient than the projection-based method in practice. This is because the projection-based method requires solving a three-dimensional convex optimization problem for each vertex. In contrast, the searching-based method only requires solving one-dimensional or two-dimensional LP for each vertex. For the searching-based method, most vertices can be solved by one-dimensional LP in the tested linear COPP3 instances, which significantly saves computational cost. In practice, we recommend the searching-and projection-based methods for linear and non-linear COPP3, respectively.
\end{remark}

\subsubsection{Guarantee of Feasibility and KKT Convergence} 

This section considers the feasibility and KKT guarantee for GOPP3. The remaining issue is the guarantee of local exactness condition \eqref{eq:KKT_approximate_reachable_set_condition} in Theorem \ref{thm:KKT_approximate_reachable_set} for the above two methods.

For the projection-based method in Section \ref{sub2sec:backward_set_copp3_extend}, assume that condition \eqref{eq:KKT_approximate_reachable_set_condition} holds for $\mathcal{B}_{\text{c}\,k+1}$. Since $\vx_{k+1}^{(q)}\in\mathcal{B}_{\text{c}\,k+1}$, we have $\vx_{k}^{(q)}\in\mathcal{B}_{\text{e}k}$. Therefore, the point $\vx_{k}^{(q)}$ can be merged into $\mathcal{B}_{\text{c}k}$ by computing the convex hull of $\mathcal{B}_{\text{c}k}$ and $\vx_{k}^{(q)}$. The resulting set $\mathcal{B}_{\text{c}k}'$ is still a convex subset of $\mathcal{B}_{\text{e}k}$. If $\vx_{k}^{(q)}$ lies in the interior of $\mathcal{B}_{\text{e}k}$, then $\mathcal{B}_{\text{c}k'}$ can be updated along some feasible directions, such as tangent directions, through one-dimensional optimization. Otherwise, $\vx_{k}^{(q)}$ may lie on a curve boundary of $\mathcal{B}_{\text{e}k}'$. In the latter case, $\mathcal{B}_{\text{c}k'}$ can be updated by searching several boundary points of $\mathcal{B}_{\text{e}k}$ along the curve boundary near $\vx_{k}^{(q)}$.

For the searching-based method in Section \ref{sub2sec:backward_set_copp3_clip} to compute $\mathcal{B}_{\text{c}k}$, condition \eqref{eq:KKT_approximate_reachable_set_condition} can be easily guaranteed. Specifically, when applying the ear-clipping method, triangles that intersect with $\mathbb{B}(\vx_k^{(q)},\delta)$ are retained and not clipped.

\begin{table*}[!t]
    \centering
    \caption{Performance comparison of different methods for TOPP2, GOPP2, TOPP3, and GOPP3 (mean $\pm$ std).}
    \label{tab:performance_comparison}
    \subcaptionbox{Time-Optimal Path Parameterization (TOPP)\label{tab:performance_comparison_topp}}{
    \resizebox{0.48\textwidth}{!}{
    \begin{tabular}{lllll}
        \toprule
        Problem & Method & SR (\%) & $T_\text{c}$ (ms) & $J_1$ \\
        \midrule
        \multirow{5}{*}{TOPP2} & TOPP2-RA & 100 & 0.433$\pm$0.031 & 30.399$\pm$1.887 \\
        & GOPP2-SOCP & 100 & 156.582$\pm$10.010 & 30.399$\pm$1.887 \\
        & GOPP2-SOCP${}^*$ & 100 & 61.616$\pm$1.550 & 30.649$\pm$1.892 \\
        & GOPP2-IDP & 100 & 13.029$\pm$0.362 & 30.412$\pm$1.886 \\
        & \textbf{GOPP2-RDDP} & \textbf{100} & \textbf{3.716$\pm$0.247} & \textbf{30.399$\pm$1.887} \\ \midrule
        \multirow{7}{*}{TOPP3} & TOPP3-RA & 100 & 24.299$\pm$1.401 & 30.768$\pm$1.886 \\ 
        & \textbf{TOPP3-RA-IRC} & 100 & 9.742$\pm$0.444 & 30.768$\pm$1.885 \\ 
        & GOPP3-SOCP & 100 & 345.298$\pm$22.141 & 30.760$\pm$1.886 \\
        & GOPP3-SOCP${}^*$ & 100 & 169.713$\pm$2.420 & 30.812$\pm$1.908 \\
        & GOPP3-IDP & 100 & 12262.139$\pm$850.531 & 32.003$\pm$1.766 \\
        & GOPP3-GDP & 62 & 69.471$\pm$3.349 & 32.699$\pm$1.773 \\
        & \textbf{GOPP3-RDDP} & \textbf{100} & \textbf{64.503$\pm$4.099} & \textbf{30.761$\pm$1.886} \\ \bottomrule
    \end{tabular}}
    }
    \subcaptionbox{General Optimal Path Parameterization (GOPP)\label{tab:performance_comparison_gopp}}{
    \resizebox{0.48\textwidth}{!}{
    \begin{tabular}{lllll}
        \toprule
        Problem & Method & SR (\%) & $T_\text{c}$ (ms) & $J_2$ \\
        \midrule
        \multirow{5}{*}{GOPP2} & TOPP2-RA & 100 & 0.431$\pm$0.037 & 58.239$\pm$2.531 \\
        & GOPP2-SOCP & 100 & 267.751$\pm$25.970 & 46.191$\pm$1.572 \\
        & GOPP2-SOCP${}^*$ & 100 & 68.429$\pm$2.327 & 46.911$\pm$3.057 \\
        & GOPP2-IDP & 100 & 14.292$\pm$0.525 & 46.261$\pm$1.566 \\
        & \textbf{GOPP2-RDDP} & \textbf{100} & \textbf{9.370$\pm$0.334} & \textbf{46.205$\pm$1.570} \\
        \midrule
        \multirow{7}{*}{GOPP3} & TOPP3-RA & 100 & 24.610$\pm$3.795 & 57.538$\pm$2.416 \\
        & \textbf{TOPP3-RA-IRC} & 100 & 9.813$\pm$0.641 & 57.529$\pm$2.418 \\ 
        & GOPP3-SOCP & 100 & 356.834$\pm$42.674 & 46.403$\pm$1.573 \\
        & GOPP3-SOCP${}^*$ & 100 & 175.139$\pm$3.258 & 48.562$\pm$4.006 \\
        & GOPP3-IDP & 100 & 11617.767$\pm$675.389 & 59.994$\pm$3.470 \\
        & GOPP3-GDP & 68 & 70.775$\pm$2.179 & 47.151$\pm$1.459 \\
        & \textbf{GOPP3-RDDP} & \textbf{100} & \textbf{61.935$\pm$0.857} & \textbf{47.131$\pm$1.550} \\ \bottomrule
    \end{tabular}}
    }
\end{table*}

\section{Experiments}

\subsection{Setup}

The proposed RDDP method and all baselines are implemented in Rust. Experiments are conducted on a laptop with an Intel${}^\text{\textregistered}$ Core${}^\text{TM}$ Ultra 9 Processor 275HX. Optimization problems in baselines are solved by Clarabel \cite{goulart2024clarabel}. The tested methods are as follows:
\begin{itemize}
    \item \textbf{TOPP2-RA} / \textbf{TOPP3-RA}: RA methods for TOPP2 \cite{pham2018new} or TOPP3 \cite{dio2025time}. All GOPP3 problems are convexified at the solution $x_1(s)$ of TOPP2-RA. In TOPP3-RA, the maximal number of vertices for the conservative backward reachable set is set to 16.
    \item \textbf{TOPP3-RA-IRC} (Ours): In contrast to TOPP3-RA, we apply the proposed improved reachable set (IRC) computation method in Section \ref{subsec:ardp_copp3} to TOPP3-RA-IRC, which should be viewed as our contribution.
    \item \textbf{GOPP2-SOCP} / \textbf{GOPP3-SOCP}: The well-known CO methods for GOPP2 \cite{verscheure2009time} or GOPP3 \cite{debrouwere2013time} formulated as second-order cone programs (SOCP). For GOPP3, the SOCP formulation is embedded into the SCP framework.
    \item \textbf{GOPP2-SOCP${}^*$} / \textbf{GOPP3-SOCP${}^*$}: GOPP2-SOCP or GOPP3-SOCP variants with a limit on the computational time for each optimization problem. In Section \ref{sub2sec:random_7axis_path}, the time limit is set to 50\,ms and 150\,ms for GOPP2 and GOPP3, respectively. The iteration is terminated when the time limit is reached, and the historical best solution is returned. In Clarabel, the real computation time is usually slightly longer than the time limit.
    \item \textbf{GOPP2-IDP} / \textbf{GOPP3-IDP}: Grid-based ideal DP method for GOPP2 or GOPP3 \cite{shin1986dynamic}. At each $s_k$, the state grid is discretized with 60 and 30$\times$30 points in GOPP2 and GOPP3, respectively.
    \item \textbf{GOPP3-GDP}: A DP method for GOPP3 based on a greedy beam search heuristic \cite{kaserer2019nearly}. The number of control segments is set to 100. For each control point, the number of sampled states is set to 25.
    \item \textbf{GOPP2-RDDP} / \textbf{GOPP3-RDDP} (Ours): Our proposed RDDP method. At each $s_k$, the initial number of sampled states is set to 5 and 5$\times$5 in GOPP2 and GOPP3, respectively. The backward and forward passes are repeated for 3 and 5 times in GOPP2 and GOPP3, respectively. Furthermore, backward reachable sets of OPP3 are computed by the proposed method in Section \ref{subsec:ardp_copp3}.
\end{itemize}


\subsection{Numerical Experiments}\label{subsec:numerical_experiments}

\subsubsection{Random 7-axis Paths for Performance Comparison}\label{sub2sec:random_7axis_path}

This section randomly generates 100 7-axis spline paths $\vq=\vq(s)$ with pure kinematic constraints. For each path, 20 waypoints are uniformly sampled in $[-1,1]^7$ and interpolated using a uniform quintic spline. The $s$-domain $[0,1]$ is discretized into 1,000 segments. The nominal maximal axial velocity, acceleration, and jerk are set to 1, 5, and 50, respectively. The final upper bounds are obtained by independently perturbing each nominal bound in each axis by a random factor uniformly sampled from $[0.9,1.1]$, while the corresponding lower bounds are set symmetrically. Two different objectives are considered: time objective $J_1=t_\text{f}$ and time-energy hybrid objective $J_2=t_\text{f}+0.1\int_0^{t_\text{f}}\sum_{i=1}^7 \tau_i^2 \, \mathrm{d}t$, where the unit point-mass dynamics $\vtau=\ddot{\vq}$ is applied for energy computation. We refer the readers to \cite[Section \RomanNum{4}--A1]{verscheure2009time} for the detailed SOCP formulation. Problems of orders 2 and 3 are both tested. In other words, four different problem settings are tested: TOPP2, GOPP2, TOPP3, and GOPP3. The quantitative results are summarized in Table \ref{tab:performance_comparison}. The success rate (SR) is determined by whether the algorithm returns a feasible trajectory for each testing case. For GOPP3-GDP, failed cases occur when the greedy pruning process eliminates all candidate grid states at some $s_k$, so no trajectory is returned and the algorithm reports infeasibility.

Table \ref{tab:performance_comparison} provides a clear comparison in terms of computational time ($T_\text{c}$), optimality ($J$), and success rate (SR). With consideration of the entire optimization problem \eqref{eq:op_discretize} as a whole, GOPP2-SOCP and GOPP3-SOCP provide high-quality reference objective values while requiring significantly more computational time than other methods. For the four cases, our RDDP approaches achieve objective values comparable to SOCP methods while requiring substantially less computational time. Specifically, the GOPP2-RDDP and GOPP3-RDDP are approximately 28.6 and 5.8 times faster than GOPP2-SOCP and GOPP3-SOCP, respectively.


Note that once a practical time limit is imposed, GOPP2-SOCP${}^*$ and GOPP3-SOCP${}^*$ exhibit evident optimality degradation. The proposed RDDP approaches achieve better objective values and higher computational efficiency than the corresponding SOCP${}^*$ methods.

\begin{figure}[!]
    \centering
    \includegraphics[width=0.9\columnwidth]{convergence.png}
    \caption{Relative optimality gap of GOPP2-IDP and our GOPP2-RDDP. (a) TOPP2. (b) GOPP2. The objective value of GOPP2-SOCP is used as the reference. This figure demonstrates the faster optimality convergence of the RDDP passes over uniform grid densification.}
    \label{fig:Convergence}
\end{figure}

\begin{figure*}[!]
    \centering
    \includegraphics[width=0.91\textwidth]{trajectories_random.png}
    \caption{Representative normalized profiles for one instance. The upper bounds are normalized to $\pm1$. (a) TOPP2. (b) GOPP2. (c) TOPP3. (d) GOPP3.}
    \label{fig:trajectories_random}
\end{figure*}


For TOPP2 and TOPP3, RA-based methods achieve the highest computational efficiency, while their greedy forward pass limits their optimality in GOPP2 and GOPP3. In TOPP2, the absolute optimality gaps of TOPP2-RA and our GOPP2-RDDP are $2.05\times10^{-4}$ and $1.99\times10^{-5}$, respectively. With fixed discretization settings, TOPP2-RA does not have a refinement loop and cannot further improve the optimality; in contrast, our GOPP2-RDDP can further reduce the optimality gap by the backward and forward passes, as shown in Fig. \ref{fig:Convergence}(a). Similarly, the optimality gap of the proposed GOPP3-RDDP is significantly smaller than that of TOPP3-RA for the TOPP3 case, especially when considering the conservatism of the backward reachable set. Furthermore, the proposed improved reachable-set computation is beneficial beyond RDDP itself. When plugged into the RA pipeline, TOPP3-RA-IRC reduces the computation time by about 60\% compared with the baseline TOPP3-RA. In summary, RA-based methods are extremely fast for pure TOPP, while the proposed RDDP methods maintain a tighter optimality gap and broader applicability to general objectives due to the objective-awareness value-function information.

In the experiment setting, the proposed RDDP methods outperform grid-based DP baselines in terms of computational time, optimality, and success rate. Compared to the IDP baseline, the proposed RDDP methods remove the restriction that trajectories must pass through predefined state-grid points, which significantly improves optimality and convergence. A comparison of the convergence between GOPP2-IDP and our GOPP2-RDDP is shown in Fig. \ref{fig:Convergence}. It can be observed that the convergence rate of  GOPP2-RDDP is much faster than that of GOPP2-IDP. Note that GOPP2-IDP improves the optimality with an increasing number of grid points, while our GOPP2-RDDP can further improve the optimality by the backward and forward passes. These results show that value-function refinement over continuous reachable domains is more effective than uniform state-grid densification for improving optimality under the tested settings.

The standard IDP method suffers from the curse of dimensionality in TOPP3 and GOPP3, which leads to a significant increase in computational time. With a greedy beam search heuristic, GOPP3-GDP significantly improves efficiency while sacrificing optimality and success rate. Note that GOPP3-GDP lacks a theoretical guarantee of feasibility, and the feasibility is empirically sensitive to the setting of the number of control segments and sampled states. In contrast, our GOPP3-RDDP maintains a $100\%$ success rate and achieves better objective values than GOPP3-GDP within less computational time.


\begin{figure*}[!]
    \centering
    \includegraphics[width=0.91\textwidth]{dim_scaling.png}
    \caption{Scalability of different methods to $n$-axis paths. (a) Computational time in TOPP problems. (b) Relative optimality gap in TOPP problems. (c) Computational time in GOPP problems. (d) Relative optimality gap in GOPP problems. In (b) and (d), the objective value of the SOCP methods is considered as the reference.}
    \label{fig:dim_scaling}
\end{figure*}

The profiles generated by the proposed RDDP methods w.r.t. the same path and constraints are shown in Fig. \ref{fig:trajectories_random}. For TOPP2 and TOPP3, the profiles satisfy the bang-singular-bang control law, which is consistent with the well-known PMP \cite{pontryagin1987mathematical}. The TOPP3 profiles take a similar shape to the TOPP2 profiles, while the acceleration are smoother due to the jerk constraints. For GOPP2 and GOPP3, the terminal time is longer than that of TOPP2 and TOPP3, respectively, while the acceleration is smaller due to the time-energy hybrid objective. These results indicate that the proposed RDDP methods can generate trajectories with different optimality and profiles by simply changing the objective function, which demonstrates the flexibility of our method.

\begin{figure*}[!]
    \centering
    \includegraphics[width=\textwidth]{real_trajecotries.png}
    \caption{Real-world experiments. (a) The robotic platform. (b) The normalized command and feedback signals in TOPP. (c) The normalized command and feedback signals in GOPP. Both TOPP and GOPP are solved by the proposed RDDP approach.}
    \label{fig:real_trajecotries}
\end{figure*}

\subsubsection{Scalability to $n$-axis Paths}

To further examine the scalability of the proposed RDDP methods, this section randomly generates 10 $n$-axis spline paths $\vq=\vq(s)$ with pure kinematic constraints, where $n=5,10,\dots,100$. The path and constraint settings are the same those in Section \ref{sub2sec:random_7axis_path}. For OPP2 and OPP3, the numbers of inequality constraints are $2n+1$ and $4n+1$ at each $s_k$, respectively.

The results are summarized in Fig. \ref{fig:dim_scaling}. As the number of axes increases, the computational time of all methods increases. The increase rate of the proposed RDDP methods is much smaller than that of the SOCP methods, which is attributed to the efficient low-dimensional DP decomposition of RDDP. Specifically, the computational time of GOPP3-RDDP becomes comparable to that of TOPP3-RA when $n$ is large, which indicates the efficiency of our searching-based method for computing the backward reachable set. For the optimality, the relative optimality gap of the proposed RDDP methods remains small as $n$ increases, which indicates the scalability of our method to high-dimensional problems. The relation of the computational time and optimality gap to $n$ is consistent with the analysis in Section \ref{sub2sec:random_7axis_path}. Therefore, the proposed RDDP methods can be efficiently scaled to high-dimensional problems while maintaining a small optimality gap.

\subsection{Real-World Experiments}\label{subsec:real_world_experiment}

Given the same OPP problem, the solutions of SOCP baselines and the proposed RDDP are highly similar. Therefore, this section conducts a real-world case study to illustrate the practical relevance of general-objective OPP, rather than to re-compare different OPP solvers. As shown in Fig. \ref{fig:real_trajecotries}(a), a 7-axis Franka Emika Panda manipulator is utilized. A quintic spline path with 16 waypoints are followed and discretized into 100 $s$-intervals. The axial constraints are set as follows: the maximal velocity $\dot{\vq}^+$ is [1.1, 1.1, 2, 1.1, 1.3, 1.3, 1.3] rad/s; the maximal acceleration $\ddot{\vq}^+$ is [6, 6, 8, 10, 12, 16, 16] rad/s${}^\text{2}$; the maximal jerk $\dddot{\vq}^+$ is [350, 490, 560, 560, 560, 560, 560] rad/s${}^\text{3}$; the maximal torque $\vtau^+$ is [52.2, 52.2, 52.2, 52.2, 7.2, 7.2, 7.2] N$\cdot$m. Note that the objective in GOPP is a hybrid of time and thermal energy. The closed-loop force control, as well as acquiring and filtering feedback signals, are implemented based on Franka-Rust \cite{Jizhou2025FrankaRust}, where trajectories planned by the proposed RDDP serve as command signals.

\begin{table}[!t]
    \centering
    \caption{Metrics for Real-World Feedback Signals.}
    \label{tab:real_world}
    \begin{tabular}{lll}
        \toprule
        Metric & TOPP3 & GOPP3 \\ \midrule
        Traversal Time $J_\text{time}$ (s) & 13.462 & 13.482 \\
        Normalized Thermal Energy $J_\text{thermal}$ (s) & 11.810 & 11.778 \\
        Normalized Total Variation of Acceleration & 720.614 & 88.862 \\ \bottomrule
    \end{tabular}
\end{table}

A time objective $J_{\text{time}}=\int_{0}^{t_\f}\mathrm{d}t$ and a time-thermal hybrid objective $J_{\text{time}}+5J_{\text{thermal}}=\int_{0}^{t_\f}\Big(1+5\norm[2]{\frac{\vtau(t)}{\vtau^+}}^2\Big)\mathrm{d}t$ are utilized for TOPP and GOPP, respectively. Some typical profiles of the proposed RDDP are shown in Fig. \ref{fig:real_trajecotries}(b)-(c), whereas quantitative results are provided in Table \ref{tab:real_world}. In this task, the gravitational component dominates the joint torques. As a result, the thermal penalty in GOPP does not substantially reduce the overall motion velocity or acceleration, where the resulting traversal time is only 0.15\% longer than that of TOPP. For the same reason, the thermal energy of the TOPP profile is also close to that of the GOPP profile.

Nevertheless, the trajectory smoothness is significantly improved by GOPP, as compared in Fig. \ref{fig:real_trajecotries}(b)-(c). Note that command signals exhibit high-frequency switching in TOPP, which lead to undesirable oscillations in feedback signals. In GOPP, the total variation of the feedback axial acceleration, normalized by $\ddot{\vq}^{+}$, is reduced by 87.7\% compared with TOPP. It is noteworthy that the total variation itself is not explicitly penalized in the GOPP objective. A potential theoretical analysis is provided as follows. A potential explanation is that $J_{\mathrm{thermal}}$ introduces strong convexity w.r.t. $x_2=\ddot{s}$, which discourages oscillatory variations in $x_2$ due to Jensen's inequality. In contrast, the jerk input in TOPP tends to saturate within its admissible bounds due to PMP. As shown in Fig. \ref{fig:real_trajecotries}(b), the planned jerk remains feasible in TOPP, while it empirically alternates between adjacent discretization intervals. The above results are consistent with Fig. \ref{fig:trajectories_random} in numerical tests.

\section{Conclusion and Future Work}

This paper has proposed reachability-augmented dual dynamic programming (RDDP) for solving optimal path parameterization problems (OPP) with general objectives and constraints. Instead of solving the large-scale optimization problem as a whole, RDDP decomposes the problem into a sequence of tractable subproblems by computing backward reachable sets. By explicitly enforcing reachability rather than relying on the relatively complete recourse assumption, the global optimality and the KKT convergence of RDDP are theoretically guaranteed under OPP-specific conditions for convex and non-convex OPP, respectively. Efficient instantiations of RDDP for second- (OPP2) and third-order OPP (OPP3) have been developed. Numerical experiments have demonstrated that RDDP simultaneously achieves high-quality objective values, feasibility, and computational efficiency across the tested OPP2 and OPP3 settings. Specifically, RDDP achieves comparable objective values to convex-optimization methods while improving computational efficiency by 28.6 and 5.8 times for OPP2 and OPP3, respectively. Compared with RA and grid-based DP baselines, RDDP better combines reachability-based feasibility preservation with objective-aware value-function guidance, which leads to high-quality feasible trajectories under general objectives. Moreover, the scalability of our RDDP methods to high-dimensional problems has been verified up to 100-axis paths. In summary, RDDP addresses a key capability gap in OPP by unifying certifiable general-objective optimization, reachability-based feasibility preservation, and online-compatible low-dimensional DP computation in a single state-grid-free framework for both OPP2 and OPP3.

In future work, the sampling strategy for cut-generation points and the conservatism of the backward reachable set can be further guided by the cut-based approximation itself. The one-step DDP optimality gap in \eqref{eq:dual_gap_k} can guide adaptive cut generation to accelerate convergence. Some advanced DDP techniques, such as cut pruning and cut reuse, can also be introduced to RDDP to further improve efficiency.

\section*{Acknowledgments}

The authors would like to thank Dr. Suqin He and Dr. Shize Lin for their expertise in TOPP.

\appendices

\section{Applicability of OPP Formulation \eqref{eq:op_discretize}}\label{app:applicability}

This appendix provides the construction of the OPP formulation \eqref{eq:op_discretize} and discusses its applicability to various OPP scenarios. Consider the description of OPP in the beginning of Section \ref{subsec:formulation_opp}. In the time domain, the state and control variables of $m$-th order OPP problems are $\hat{\vx}(t)=(s(t), \dot{s}(t), \ldots, s^{(m-1)}(t))\in\R^m$ and $\hat{u}(t)=s^{(m)}(t)\in\R$, respectively. Then, the $m$-th order OPP problem can be formulated as the following optimal control problem:
\begin{subequations}\label{eq:ocp_time}
    \begin{align}
        \min_{\hat{\vx}(\bullet), \hat{u}(\bullet), t_\text{f}} \quad &J = \int_0^{t_\text{f}} \widehat{L}(\hat{\vx}(t), \hat{u}(t)) \mathrm{d}t + \widehat\Phi(\hat{\vx}(t_\text{f})) \label{eq:ocp_time_objective}\\
        \st \quad & \dot{s}^{(i-1)}(t) = s^{(i)}(t), \,\,i=1,2,\ldots,m, \\
        &\dot{s}(t) \geq 0,\label{eq:ocp_time_dynamics}\\
        &\hat{\vf}(\hat{\vx}(t); s(t)) \leq \vzero,  \label{eq:ocp_time_state_constraints}\\
        &\hat{\vh}(\hat{\vx}(t), \hat{u}(t); s(t)) \leq \vzero, \label{eq:ocp_time_mix_constraints} \\
        &\hat{\vx}(0)\in\widehat{\mathcal{X}}_0, \,\, \hat{\vx}(t_\text{f})\in\widehat{\mathcal{X}}_\text{f}.\label{eq:ocp_time_boundary}
    \end{align}
\end{subequations}
In OPP \eqref{eq:ocp_time}, Eqs. \eqref{eq:ocp_time_dynamics}--\eqref{eq:ocp_time_mix_constraints} are set for $t\in(0,t_\text{f})$.

Following \cite{verscheure2009time,debrouwere2013time}, OPP \eqref{eq:ocp_time} can be transformed into equivalent problems in the $s$-domain to get rid of the free terminal time $t_\text{f}$ and non-convexity caused by the time variable. Denote the derivative of a variable $\bullet$ with respect to $s$ as $\bullet'=\frac{\mathrm{d}\bullet}{\mathrm{d}s}$. Let $x_1(s)=\frac12\dot{s}^2$ and $x_{i+1}(s)=x_i'(s)$ for $i=1,2,\ldots,m-1$. Then, the state and control variables in the $s$-domain are $\vx(s)=(x_i)_{i=1}^{m-1}\in\R^{m-1}$ and $u(s)=x_{m}(s)\in\R$, respectively. The $m$-th order OPP problem can be represented as the following optimal control problem in the $s$-domain:
\begin{subequations}\label{eq:ocp_parameter}
    \begin{align}
        \min_{\vx(\bullet), u(\bullet)} \quad &J = \int_0^{s_\text{f}} L(\vx(s), u(s); s) \mathrm{d}s + \Phi(\vx(s_\text{f})) \label{eq:ocp_parameter_objective}\\
        \st \quad & x_i'(s) = x_{i+1}(s), \,\,i=1,2,\ldots,m-1, \\
        &\vf(\vx(s); s) \leq \vzero,  \label{eq:ocp_parameter_state_constraints}\\
        &\vh(\vx(s), u(s); s) \leq \vzero, \label{eq:ocp_parameter_mix_constraints} \\
        &\vx(0)\in \mathcal{X}_0, \,\, \vx(s_\text{f})\in \mathcal{X}_\text{f}.\label{eq:ocp_parameter_boundary}
    \end{align}
\end{subequations}
Specifically, the constraint $x_1(s)\geq0$, $\forall s\in(0,s_\text{f})$ is contained in the state constraints \eqref{eq:ocp_parameter_state_constraints}.

Generally, OPP2 \cite{verscheure2009time,pham2018new} and OPP3 \cite{debrouwere2013time,wang2026online} problems are of interest. Denote
\begin{equation}\label{eq:state_control_relation}
    x_2(s)=x_1'(s)=\frac{\dot{x}_1}{\dot{s}}=\ddot{s},\,\,x_3(s)=x_2'(s)=\frac{\dot{x}_2}{\dot{s}}=\frac{\dddot{s}}{\dot{s}}.
\end{equation}
For OPP2, the decision variables are $\vx(s)=x_1(s)$ and $u(s)=x_2(s)$, respectively. For OPP3, the decision variables are $\vx(s)=(x_1(s), x_2(s))$ and $u(s)=x_3(s)$, respectively.

It has been recognized that OPP \eqref{eq:ocp_parameter} provides a highly expressive framework, capable of modeling various objectives and constraints. Some examples are provided as follows.

\begin{example}
    Consider the axial torque constraint $\underline{\vtau}\leq\vtau(s)\leq \overline\vtau$. Denote the dynamics of the robotic system as
    \begin{equation}
        \mM(\vq)\ddot{\vq} + \mC(\vq, \dot{\vq})\dot{\vq} + \vg(\vq) = \vtau,
    \end{equation}
    where $\mM(\vq)$, $\mC(\vq, \dot{\vq})$, and $\vg(\vq)$ are the inertia matrix, the Coriolis and centrifugal matrix, and the gravity vector, respectively. Specifically, $\mC(\vq, \dot{\vq})$ is linear w.r.t. $\dot{\vq}$. Note that
    \begin{subequations}
        \begin{align}
            &\dot{\vq}=\vq'\dot{s}=\vq'\sqrt{2x_1},\label{eq:velocity}\\
            &\ddot{\vq}=\vq''\dot{s}^2+\vq'\ddot{s}=2\vq''x_1+\vq'x_2.\label{eq:acceleration}
        \end{align}
    \end{subequations}
    Then, the torque can be expressed as
    \begin{equation}\label{eq:torque}
        \vtau(s) = \mM(\vq)(2\vq''x_1+\vq'x_2) + 2\mC(\vq, \vq')\vq'x_1 + \vg(\vq),
    \end{equation}
    which is affine in $x_1$ and $x_2$. In OPP2, the torque constraint can be represented as the mixed constraints \eqref{eq:ocp_parameter_mix_constraints}. In OPP3, it can be represented as the state constraints \eqref{eq:ocp_parameter_state_constraints}.

    Similar analysis can be applied to the box constraints on axial velocity \eqref{eq:velocity} and acceleration \eqref{eq:acceleration}, both of which can be transformed into linear inequalities in $x_1$ and $x_2$.
\end{example}

\begin{example}
    For OPP3, the constraint on the axial jerk $\dddot{\vq}$ can be represented as the mixed constraints \eqref{eq:ocp_parameter_mix_constraints}. Since
    \begin{equation}
        \dddot{\vq}=\vq'''\dot{s}^3 + 3\vq''\dot{s}\ddot{s} + \vq'\dddot{s} = \sqrt{2x_1}(2\vq'''x_1 + 3\vq''x_2 + \vq'x_3),
    \end{equation}
    the jerk constraint $\underline{\dddot{\vq}}\leq \dddot{\vq}\leq \overline{\dddot{\vq}}$ can be represented as
    \begin{equation}\label{eq:jerk_constraint}
        \frac{1}{\sqrt{2x_1}}\underline{\dddot{\vq}} \leq 2\vq'''x_1 + 3\vq''x_2 + \vq'x_3 \leq \frac{1}{\sqrt{2x_1}}\overline{\dddot{\vq}}.
    \end{equation}
    In \eqref{eq:jerk_constraint}, the limits of jerk usually satisfy $\underline{\dddot{\vq}}<\vzero<\overline{\dddot{\vq}}$. Then, the jerk constraint is concave in $x_1$. Similar analysis can be applied to the constraint on the axial torque rate $\dot{\vtau}$ \cite{debrouwere2013time,wang2026online}.
\end{example}

\begin{example}
    \begin{subequations}
        Consider some common objectives. In TOPP, the objective is to minimize the terminal time
        \begin{equation}\label{eq:objective_time}
            J_1=\int_0^{t_\text{f}} \mathrm{d}t = \int_0^{s_\text{f}} \frac{1}{\dot{s}} \mathrm{d}s = \int_0^{s_\text{f}} \frac{1}{\sqrt{2x_1(s)}} \mathrm{d}s.
        \end{equation}
        \cite{verscheure2009time} further considers the thermal energy
        \begin{equation}\label{eq:objective_thermal}
            J_2 = \int_0^{t_\text{f}} \vtau^\top \mQ\vtau \mathrm{d}t = \int_0^{s_\text{f}} \frac{\vtau^\top\mQ\vtau}{\sqrt{2x_1}} \mathrm{d}s,\,\mQ\succeq \vzero,
        \end{equation}
        and the total variation of the torque
        \begin{equation}\label{eq:objective_torque_variation}
            J_3 = \int_0^{t_\text{f}} \norm[]{\dot{\vtau}} \mathrm{d}t= \int_0^{s_\text{f}} \frac{\norm[]{\vtau'\dot{s}}}{\dot{s}} \mathrm{d}s = \int_0^{s_\text{f}} \norm[]{\vtau'} \mathrm{d}s.
        \end{equation}
    \end{subequations}
    In OPP2, the linear combination of \eqref{eq:objective_time} and \eqref{eq:objective_thermal} can be represented as the objective \eqref{eq:ocp_parameter_objective}. In OPP3, the linear combination of all the three objectives can be handled.
\end{example}

Finally, OPP \eqref{eq:ocp_parameter} can be discretized into the formulation \eqref{eq:op_discretize}. All inequality constraints are enforced at the grid points $s_k$ or the two-side limits $s_k^\pm$. Specifically, the state constraints \eqref{eq:op_discretize_state_constraints} refer to $\vf(\vx(s_k); s_k)=\vf(\vx_k; s_k)\leq\vzero$ in \eqref{eq:ocp_parameter_state_constraints}, whereas the mixed constraints \eqref{eq:op_discretize_mix_constraints} refer to $\vh(\vx(s_k), u(s_k^+); s_k)=\vh(\vx_k, u_k; s_k)\leq\vzero$ for $k<N$ and $\vh(\vx(s_k), u(s_k^-); s_k)=\vh(\vx_k, u_{k-1}; s_k)=\vh(\vA_{k}\vx_{k-1}+\vb_{k-1}u_{k-1}, u_{k-1}; s_k)\leq\vzero$ for $k>0$. Furthermore, the objective function $L_k$ in \eqref{eq:op_discretize_objective} is a numerical approximation of the integral $\int_{s_{k-1}}^{s_k} L(\vx(s), u(s); s) \mathrm{d}s$ in \eqref{eq:ocp_parameter_objective}. In summary, all the functions and sets in \eqref{eq:op_discretize} are given, whereas $\vx_k$ and $u_k$ for all $k$ are the decision variables.

\section{Proofs of Theorems and Propositions}\label{app:proofs}

\begin{lemma}\label{lemma:min_is_convex}
    Consider a convex function $V:\Omega\subset\R^m\times\R^n\to\R$ where $\Omega$ is a non-empty convex set. Define the function $U(\vx)=\min_{\vy\in\mathcal{Y}(\vx)} V(\vx,\vy)$, where the domain of $U$ is $\dom U=\{\vx\in\R^m\mid \mathcal{Y}(\vx)\neq\varnothing\}$, and $\mathcal{Y}(\vx)=\{\vy\in\R^n\mid (\vx,\vy)\in\Omega\}$. Then, $U$ is convex in the convex set $\dom U$.
\end{lemma}

\begin{proof}
    Denote the epigraph of $V$ and $U$ as $\text{epi}(V)$ and $\text{epi}(U)$, respectively. By the definition of $U$, $\text{epi}(U)$ is the projection of $\text{epi}(V)$ onto the $\vx$-subspace, which is convex since $\text{epi}(V)$ is convex. Therefore, $U$ is convex in $\dom U$.
\end{proof}

\begin{proof}[Proof of Lemma \ref{lemma:convexity_reachset_valfunc}]
    Consider the following convex set
    \begin{align}
        \widehat{\mathcal{B}}_k=\Big\{&(\vx_{k}, \ldots, \vx_N, u_k, \ldots, u_{N-1})\mid \notag\\
        &\forall k\leq j<N,\,\vx_{j+1} = \vA_j\vx_j+\vb_j u_j, \notag\\
        &\vf_j(\vx_j) \leq \vzero,\,\vh_j(\vx_j, u_j) \leq \vzero,\,\vx_N \in \mathcal{X}_\text{f} \Big\}.\label{eq:widehat_backward_set}
    \end{align}
    By Assumption \ref{assum:feasibility_smoothness}, $\widehat{\mathcal{B}}_k$ is compact. As a projection of $\widehat{\mathcal{B}}_k$ onto the $\vx_k$-subspace, $\mathcal{B}_k$ is also convex and compact. By Lemma \ref{lemma:min_is_convex}, $V_k^*$ is convex in $\mathcal{B}_k$ since $\widehat{V}_k:\widehat{\mathcal{B}}_k\to\R$ is convex:
    \begin{equation}
        \widehat{V}_k\left((\vx_j)_{j=k}^N, (u_j)_{j=k}^{N-1}\right) = \sum_{j=k}^{N-1} L_j(\vx_j, u_j) + \Phi(\vx_N).
    \end{equation}
    The convexity and compactness of $\mathcal{C}_k(\vx_k)$ is straightforward.

    If all the constraints are linear inequalities, then $\widehat{\mathcal{B}}_k$ is a polytope, and thus $\mathcal{B}_k$ is also a polytope.
\end{proof}

\begin{proof}[Proof of Proposition \ref{prop:Lipschitz}]
    We prove this proposition by induction. For $k=N$, \eqref{eq:approximation_terminal_value} and Assumption \ref{assum:Lipschitz} imply that $\mathcal{V}_N^p=\Phi$ is convex and locally Lipschitz continuous in $\mathcal{B}_N=\mathcal{X}_\text{f}$. Then, $\forall \vx_N\in\mathcal{B}_N$, $\exists \varepsilon_N^V(\vx_N)>0$ small enough and $C_N^V(\vx_N)>0$ independent of $p$, s.t. $\forall \vy_1,\vy_2\in\mathcal{B}_N\cap\mathbb{B}(\vx_N,\varepsilon_N^V(\vx_N))$, we have $\abs{\mathcal{V}_N^p(\vy_1)-\mathcal{V}_N^p(\vy_2)}\leq C_N^V(\vx_N)\norm[2]{\vy_1-\vy_2}$. Furthermore, the Clarke subgradient of $\mathcal{V}_N^p=\Phi$ exists and is controlled by $C_N^V$ due to its local Lipschitz continuity \cite{clarke1990optimization}.

    For $k<N$, assume that the conclusion holds for $\mathcal{V}_{k+1}^p$. Fix $\vx\in\mathcal{B}_k$. We aim to prove that the conclusion also holds for $\widehat{\mathcal{V}}_k^p$ and $\mathcal{V}_k^p$. Let $\widehat{J}_k(u;\vy)=L_k(\vy, u) + \mathcal{V}_{k+1}^p(\vA_k\vy+\vb_k u)$. By \eqref{eq:approximate_value_function_backward}, we have $\forall \vy\in\mathcal{B}_k$,
    \begin{subequations}\label{eq:approximation_proof}
        \begin{align}
            \widehat{\mathcal{V}}_k^p(\vy) = \min\,\,& \widehat{J}_k(u;\vy) \\
            \st\,\,&\underbar{u}_k(\vy)\leq u\leq \bar{u}_k(\vy).
        \end{align}
    \end{subequations}
    According to Lemma \ref{lemma:min_is_convex}, $\widehat{\mathcal{V}}_k^p$ is convex. In the following, we only prove the local Lipschitz continuity of $\widehat{\mathcal{V}}_k^p$, whereas the existence as well as the boundedness of its Clarke subgradient is straightforward as a corollary \cite{clarke1990optimization}.

    By the induction hypothesis, $\forall u\in\mathcal{C}_k(\vx)$, $\exists \varepsilon_J(\vx,u)>0$, $C_{\vx}^J(\vx,u)>0$, and $C_{u}^J(\vx,u)>0$, s.t. $\forall \vy_1,\vy_2\in\mathcal{B}_k\cap\mathbb{B}(\vx,\varepsilon_J(\vx,u))$, $v_1\in\mathcal{C}_k(\vy_1)\cap\mathbb{B}(u,\varepsilon_J(\vx,u))$, and $v_2\in\mathcal{C}_k(\vy_2)\cap\mathbb{B}(u,\varepsilon_J(\vx,u))$, we have $\abs{\widehat{J}_k(v_1;\vy_1)-\widehat{J}_k(v_2;\vy_2)}\leq C_{\vx}^J(\vx,u)\norm[2]{\vy_1-\vy_2}+C_{u}^J(\vx,u)\abs{v_1-v_2}$. Note that $\varepsilon_J$, $C_{\vx}^J$, and $C_{u}^J$ can be constructed by the local Lipschitz constants of $L_k$ and $\mathcal{V}_{k+1}^p$, which are all independent of $p$.

    By Assumption \ref{assum:Lipschitz}, $\underbar{u}_k$ and $\bar{u}_k$ are both locally Lipschitz continuous. Therefore, $\exists \varepsilon_u(\vx)>0$, $C_k^u(\vx)>0$, s.t. $\forall \vy_1,\vy_2\in\mathcal{B}_k\cap\mathbb{B}(\vx,\varepsilon_u(\vx))$, we have $\abs{\underbar{u}_k(\vy_1)-\underbar{u}_k(\vy_2)}$, $\abs{\bar{u}_k(\vy_1)-\bar{u}_k(\vy_2)}\leq C_k^u(\vx)\norm[2]{\vy_1-\vy_2}$. The constants $\varepsilon_u(\vx)$ and $C_k^u(\vx)$ are also independent of $p$.

    Consider the open cover of the compact set $\mathcal{C}_k(\vx)\subset\bigcup_{u\in\mathcal{C}_k(\vx)}\mathbb{B}(u,\varepsilon_J(\vx,u))$. Select a finite subcover $\mathcal{C}_k(\vx)\subset\bigcup_{i=1}^M\mathbb{B}(u^i,\varepsilon_J(\vx,u^i))$ where $\bar{u}_k(\vx)$ and $\underbar{u}_k(\vx)$ are both contained in $\{u^i\}_{i=1}^M$. Define two positive constants:
    \begin{subequations}
        \begin{align}
            &\varepsilon_k^V(\vx)=\min\left\{\varepsilon_u(\vx), \min_{1\leq i\leq M}\frac{\varepsilon_J(\vx,u^i)}{1+C_k^u(\vx)}\right\},\\
            &C_k^V(\vx)=\max_{1\leq i\leq M}\Big\{C_{\vx}^J(\vx,u^i)+C_{u}^J(\vx,u^i)C_k^u(\vx)\Big\},
        \end{align}
    \end{subequations}
    both of which are independent of $p$. 

    Consider any $\vy_1,\vy_2\in\mathcal{B}_k\cap\mathbb{B}(\vx,\varepsilon_k^V(\vx))$. Select a one-step optimal control $v^*\in\argmin_{u\in\mathcal{C}_k(\vy_1)} \widehat{J}_k(u;\vy_1)$. Note that $\mathcal{C}_k(\vy_1)\subset\mathbb{B}(\mathcal{C}_k(\vx),C_k^u(\vx)\varepsilon_k^V(\vx))\subset\bigcup_{i=1}^M\mathbb{B}(u^i,\varepsilon_J(\vx,u^i))$ due to the fact that $\bar{u}_k(\vx),\underbar{u}_k(\vx)\in\{u^i\}_{i=1}^M$. Denote a feasible control $\hat{v}=\max\{\min\{u^*,\bar{u}_k(\vy_2)\},\underbar{u}_k(\vy_2)\}\in\mathcal{C}_k(\vy_2)$. Then, we construct some convex combinations of $(\vy_1,v^*)$ and $(\vy_2,\hat{v})$ as middle points. Specifically, let $0=\lambda_0<\lambda_1<\cdots<\lambda_I=1$. For each $1\leq i\leq I$, define $\vy^i=\lambda_i\vy_1+(1-\lambda_i)\vy_2$ and $v^i=\lambda_i v^*+(1-\lambda_i)\hat{v}$. Without loss of generality, assume that $v^i\in\mathbb{B}(u^{j_i},\varepsilon_J(\vx,u^{j_i}))\cap\mathbb{B}(u^{j_{i-1}},\varepsilon_J(\vx,u^{j_{i-1}}))$ for some $1\leq j_i\leq M$. Then, we have
    \begin{align}
        &\widehat{\mathcal{V}}_k^p(\vy_2)-\widehat{\mathcal{V}}_k^p(\vy_1)\notag\\
        \leq& \widehat{J}_k(\hat{v};\vy_2) - \widehat{J}_k(v^*;\vy_1)\notag\\
        =&\sum_{i=1}^{I}\left[\widehat{J}_k(v^i;\vy^i) - \widehat{J}_k(v^{i-1};\vy^{i-1})\right]\notag\\
        \leq& \sum_{i=1}^{I}\left[C_{\vx}^J(\vx,u^{j_i})\norm[2]{\vy^i-\vy^{i-1}}+C_{u}^J(\vx,u^{j_i})\abs{v^i-v^{i-1}}\right]\notag\\
        \leq& \sum_{i=1}^{I}\left[(C_{\vx}^J(\vx,u^{j_i})+C_{u}^J(\vx,u^{j_i})C_k^u(\vx))\norm[2]{\vy^i-\vy^{i-1}}\right]\notag\\
        \leq& C_k^V(\vx)\sum_{i=1}^{I}\norm[2]{\vy^i-\vy^{i-1}}
        =C_k^V(\vx)\norm[2]{\vy_2-\vy_1}.
    \end{align}
    For the same reason, we also have $\widehat{\mathcal{V}}_k^p(\vy_1)-\widehat{\mathcal{V}}_k^p(\vy_2)\leq C_k^V(\vx)\norm[2]{\vy_2-\vy_1}$. As a result, we have
    \begin{equation}\label{eq:proof_Lipschitz_V}
        \abs{\widehat{\mathcal{V}}_k^p(\vy_2)-\widehat{\mathcal{V}}_k^p(\vy_1)}\leq C_k^V(\vx)\norm[2]{\vy_2-\vy_1},
    \end{equation}
    which provides the local Lipschitz continuity of $\widehat{\mathcal{V}}_k^p$ at $\vx$.

    According to \cite{clarke1990optimization}, $C_k^V(\vx)$ also controls the Clarke subgradient of $\widehat{\mathcal{V}}_k^p$ at $\vx$, $\forall \vx\in\mathcal{B}_k$. By \eqref{eq:support_plane}, $\mathcal{V}_k^p$ is a maximum of finitely many supporting hyperplanes of $\widehat{\mathcal{V}}_k^p$; hence, $\mathcal{V}_k^p$ is also convex and locally Lipschitz continuous with the same local Lipschitz constant $C_k^V$.

    Note that the above analysis is only dependent on the local Lipschitz continuity of $L_k$, $\Phi$, $\underbar{u}_k$, and $\bar{u}_k$. For the same reason, the conclusion with the same local Lipschitz constant $C_k^V$ also holds for the real value function $V_k^*$. Note that $C_k^V$ is independent of $p$.
\end{proof}

\begin{proof}[Proof of Proposition \ref{prop:bound_of_value_function}]
    According to Proposition \ref{prop:Lipschitz}, $\widehat{\mathcal{V}}_k^p$ is convex and has Clarke subgradient. In the definition \eqref{eq:support_plane}, for each $i$, we have $\widehat{\mathcal{V}}_k^p(\vx_k)\geq\widehat{\mathcal{V}}_k^p(\vx_k^i)+{\vg_k^i}^\top(\vx_k-\vx_k^i)$. Therefore, $\widehat{\mathcal{V}}_k^p(\vx_k)\geq\mathcal{V}_k^p(\vx_k)$ holds. According to the definition of the real value function \eqref{eq:def_value_function}, $V_k^*(\vx_k)\leq J_k(\bx_k,\bu_k)$ also holds.

    Next, we prove that $\widehat{\mathcal{V}}_k^p(\vx_k)\leq V_k^*(\vx_k)$ by induction. For $k=N-1$, $\mathcal{V}_N^p(\vx_N)=V_N^*(\vx_N)=J_N(\bx_N,\bu_N)$ holds. By comparing \eqref{eq:def_value_function} and \eqref{eq:approximate_value_function_backward}, we have $\widehat{\mathcal{V}}_{N-1}^p(\vx_{N-1})=V_{N-1}^*(\vx_{N-1})$. For $k<N-1$, assume that $\widehat{\mathcal{V}}_{k+1}^p(\vx_{k+1})\leq V_{k+1}^*(\vx_{k+1})$ holds for each $\vx_{k+1}\in\mathcal{B}_{k+1}$. For any $\vx_k\in\mathcal{B}_k$, we have $\mathcal{V}_{k}^p(\vx_{k})\leq V_{k}^*(\vx_{k})$ since $\mathcal{V}_{k+1}^p(\vx_{k+1})\leq\widehat{\mathcal{V}}_{k+1}^p(\vx_{k+1})\leq V_{k+1}^*(\vx_{k+1})$. By comparing \eqref{eq:def_value_function} and \eqref{eq:approximate_value_function_backward}, we have $\widehat{\mathcal{V}}_k^p(\vx_k)\leq V_k^*(\vx_k)$.
\end{proof}

\begin{proof}[Proof of Theorem \ref{thm:near_optimality_gap}]
    Evidently, $J_0(\bx_0^{p*},\bu_0^{p*}) \geq J^*$ holds. Arbitrarily select a globally optimal solution $\bx_0^*=(\vx_k^*)_{k=0}^N$ and $\bu_0^*=(u_k^*)_{k=0}^{N-1}$ of COPP \eqref{eq:op_discretize}. By \eqref{eq:refine_x0}, we have $\mathcal{V}_0^p(\vx_0^{p*}) \leq \mathcal{V}_0^p(\vx_0^*)$. By Proposition \ref{prop:bound_of_value_function}, we have $\mathcal{V}_0^p(\vx_0^{*}) \leq V_0^*(\vx_0^{*}) = J^*$. Hence, we have $\varepsilon_\text{gap}(\bx_0^{p*},\bu_0^{p*}) = J_0(\bx_0^{p*},\bu_0^{p*}) - \mathcal{V}_0^p(\vx_0^{p*}) \geq J_0(\bx_0^{p*},\bu_0^{p*}) - \mathcal{V}_0^p(\vx_0^{*}) \geq J_0(\bx_0^{p*},\bu_0^{p*}) - J^* \geq 0$.
\end{proof}

\begin{proof}[Proof of Theorem \ref{thm:adrp_optimal}]
    Consider the solution sequence $\{(\bx_0^{p*},\bu_0^{p*})\}_{p=0}^\infty$ generated by RDDP. According to Assumption \ref{assum:feasibility_smoothness}, the feasible set of COPP \eqref{eq:op_discretize}, denoted by $\mathcal{F}$, is compact. We select a convergent subsequence $\{(\bx_0^{p_i*},\bu_0^{p_i*})\}_{i=0}^\infty$ with the feasible limit point $(\bx_0^*,\bu_0^*)\in\mathcal{F}$, where $\bx_0^*=(\vx_k^*)_{k=0}^N$ and $\bu_0^*=(u_k^*)_{k=0}^{N-1}$. Next, we prove that $(\bx_0^*,\bu_0^*)$ is a globally optimal solution of COPP \eqref{eq:op_discretize}, which implies \eqref{eq:convergence_condition} due to the local Lipschitz continuity of the real cost function $J_0$.

    According to Proposition \ref{prop:Lipschitz}, for each $k$, convex functions $\widehat{\mathcal{V}}_k^p$ and $\mathcal{V}_k^p$ share the same local Lipschitz constant $C_k^V(\vx_k^*)$ and the associated neighborhood radius $\varepsilon_k^V(\vx_k^*)$ at $\vx_k^*\in\mathcal{B}_k$ across $p$. Arbitrarily select a small $\delta\in(0,\min_{k}\varepsilon_k^V(\vx_k^*))$. By the convergence of $\{\bx_0^{p_i*}\}_{i=0}^\infty$ to $\bx_0^*$, there exists a large index $I(\delta)\in\N^*$, s.t. $\forall i\geq I(\delta)$, we have $\norm[2]{\vx_k^{p_i*}-\vx_k^*}<\delta<\varepsilon_k^V(\vx_k^*)$ for all $k$.

    Consider any $i>j\geq I(\delta)$. According to the RDDP sampling process \eqref{eq:ardp_more_sample}, for each $k$, the state $\vx_k^{p_j*}\in\mathcal{B}_k$ has been sampled to construct $\mathcal{V}_k^{p_i}$ in \eqref{eq:support_plane}. Therefore, we have 
    \begin{equation}\label{eq:proof_ardp_previous_sample}
        \widehat{\mathcal{V}}_k^{p_i}(\vx_k^{p_j*})=\mathcal{V}_k^{p_i}(\vx_k^{p_j*}).
    \end{equation}
    According to \eqref{eq:widehat_V_L_V_refine}, \eqref{eq:proof_ardp_previous_sample}, and Proposition \ref{prop:Lipschitz}, we have
    \begin{align}
        &\varepsilon_\text{gap}(\bx_0^{p_i*},\bu_0^{p_i*})\notag\\
        =&\Phi(\vx_N^{p_i*}) - \mathcal{V}_0^{p_i}(\vx_0^{p_i*})+\sum_{k=0}^{N-1} L_k(\vx_k^{p_i*}, u_k^{p_i*})\notag\\
        =&\sum_{k=0}^{N-1}\left[L_k(\vx_k^{p_i*}, u_k^{p_i*})+\mathcal{V}_{k+1}^{p_i}(\vx_{k+1}^{p_i*})-\mathcal{V}_k^{p_i}(\vx_k^{p_i*})\right]\notag\\
        =&\sum_{k=0}^{N-1}\left[\widehat{\mathcal{V}}_k^{p_i}(\vx_k^{p_i*})-\mathcal{V}_k^{p_i}(\vx_k^{p_i*})\right]\notag\\
        \leq&\sum_{k=0}^{N-1}\left[\widehat{\mathcal{V}}_k^{p_i}(\vx_k^{p_j*})-\mathcal{V}_k^{p_i}(\vx_k^{p_j*})+4C_k^V(\vx_k^*)\delta\right]\notag\\
        =&4\delta\sum_{k=0}^{N-1}C_k^V(\vx_k^*),
    \end{align}
    where the constant coefficient $4\sum_{k=0}^{N-1}C_k^V(\vx_k^*)$ is independent of $i$, $j$ and $I$. Fix $j$ and $I$, and let $i\to\infty$. Then,
    \begin{equation}
        \forall \delta>0,\,\limsup_{i\to\infty} \varepsilon_\text{gap}(\bx_0^{p_i*},\bu_0^{p_i*}) \leq 4\delta\sum_{k=0}^{N-1}C_k^V(\vx_k^*).
    \end{equation}
    The left-hand side is independent of $\delta$. Hence, we have
    \begin{align}
        0\leq&\liminf_{p\to\infty} \varepsilon_\text{gap}(\bx_0^{p*},\bu_0^{p*})\notag\\
        \leq&\limsup_{i\to\infty} \varepsilon_\text{gap}(\bx_0^{p_i*},\bu_0^{p_i*})=0,
    \end{align}
    which proves the theorem.
\end{proof}

\begin{proof}[Proof of Theorem \ref{thm:KKT_exact_reachable_set}]
    Theorem \ref{thm:KKT_exact_reachable_set} is a direct application of \cite[Theorem 10]{sriperumbudur2009convergence}. Specifically, Assumption \ref{assum:feasibility_smoothness} guarantees the consistent compactness of the feasible set across iterations, noting that the convexified feasible set is a subset of the original compact feasible set. Furthermore, Theorem \ref{thm:adrp_optimal} guarantees the global optimality of each iteration. These conditions collectively satisfy the assumptions for \cite[Theorem 10]{sriperumbudur2009convergence}.
\end{proof}

\begin{proof}[Proof of Theorem \ref{thm:KKT_approximate_reachable_set}]
    The proof is built upon Zangwill's global convergence theorem \cite[Section 4.5]{zangwill1969nonlinear}. Three conditions should be validated: (i) monotonicity of the objective function, (ii) compactness of the solution sequence, and (iii) closedness of the algorithmic mapping. Specifically, Conditions (ii) and (iii) are guaranteed by \cite{sriperumbudur2009convergence}, which remains valid for the conservative problem \eqref{eq:op_discretize_convexity_conservative}.

    For iteration $q\geq1$, denote the feasible sets of the convexified problem with and without constraints \eqref{eq:op_discretize_convexity_conservative_additional_constraints} by $\mathcal{F}_{\text{c}}^{(q)}$ and $\mathcal{F}_{\text{e}}^{(q)}$, respectively. Then, $(\bx_0^{(q)},\bu_0^{(q)}),\,(\bx_0^{(q+1)},\bu_0^{(q+1)})\in\mathcal{F}_{\text{c}}^{(q)}\subset\mathcal{F}_{\text{e}}^{(q)}$ holds. By \eqref{eq:KKT_approximate_reachable_set_condition}, there exists a constant $\delta'>0$ independent of $q$, s.t. $\mathcal{F}_{\text{c}}^{(q)}\cap\mathbb{B}(\bx_0^{(q)},\delta')=\mathcal{F}_{\text{e}}^{(q)}\cap\mathbb{B}(\bx_0^{(q)},\delta')$. Specifically, $\delta$ can be constructed by $\delta'$ for the local exactness condition. Since the objective function \eqref{eq:op_discretize_convexity_conservative_objective} is convex, Condition (i) also holds. Therefore, all conditions of Zangwill's global convergence theorem are satisfied, which guarantees the KKT convergence of the above process.
\end{proof}

\begin{proof}[Proof of Corollary \ref{cor:global_optimal_approximate_reachable_set}]
    Corollary \ref{cor:global_optimal_approximate_reachable_set} is a direct application of Theorem \ref{thm:KKT_approximate_reachable_set}, as the KKT condition is a necessary and sufficient optimality condition for convex optimization problems under certain CQ.
\end{proof}

\ifCLASSOPTIONcaptionsoff
    \newpage
\fi

\bibliographystyle{myIEEEtran}
\bibliography{IEEEabrv,refs/ref}

\end{document}
